\documentclass[a4paper,10pt]{jsarticle}
\usepackage{amsmath,amssymb,amsthm}
\usepackage[dvipdfmx]{graphicx}
\usepackage{booktabs}
\usepackage{array}
\usepackage{url}

\usepackage[margin=25mm]{geometry}

\theoremstyle{definition}
\newtheorem{definition}{定義}[section]
\newtheorem{proposition}{命題}[section]
\newtheorem{lemma}{補題}[section]
\newtheorem{theorem}{定理}[section]
\newtheorem{conjecture}{予想}[section]
\newtheorem{remark}{注意}[section]
\makeatletter
\renewenvironment{proof}[1][\proofname]{\par
  \pushQED{\qed}%
  \normalfont \topsep6\p@\@plus6\p@\relax
  \trivlist
  \item[\hskip\labelsep\textbf{#1}]\ignorespaces
}{%
  \popQED\endtrivlist\@endpefalse
}
\renewcommand{\proofname}{証明}
\makeatother

\newcommand{\CVaR}{\mathrm{CVaR}}
\newcommand{\Wone}{W_1}
\expandafter\let\csname Finv\endcsname\undefined
\newcommand{\Finv}{F^{-1}}
\newcommand{\Fhinv}{\hat{F}^{-1}}
\newcommand{\R}{\mathbb{R}}
\newcommand{\E}{\mathbb{E}}
\newcommand{\Prob}{\mathbb{P}}
\newcommand{\dist}{\mathcal{D}_g}
\newcommand{\todo}[1]{\{#1\}}

\title{\vspace{-1cm}}
\author{}
\date{}

\begin{document}

%==================== 表紙 ====================
\begin{titlepage}
\begin{center}
\vspace*{3cm}
{\LARGE 卒業研究　中間レポート}\\[1.5cm]
{\Large 歪み関数重み付き1-Wasserstein距離に基づく\\[0.3cm]
quantile fraction学習と\\[0.3cm]
CVaR推定誤差の低減}\\[0.6cm]
{\large ---分布型強化学習FQFの歪みリスク尺度への適応---}\\[3cm]

\begin{tabular}{ll}
学籍番号 & 【学籍番号】\\[0.3cm]
氏名 & 【氏名】\\[0.3cm]
所属 & 【研究室名】\\[0.3cm]
指導教員 & 【指導教員名】\\[0.3cm]
提出日 & 2026年9月14日\\
\end{tabular}
\vfill
\end{center}
\end{titlepage}

\tableofcontents
\newpage

%==================== 1 はじめに ====================
\section{はじめに}
\label{sec:intro}

\subsection{研究動機}
強化学習における価値関数 $Q^\pi(s,a)=\E[Z^\pi(s,a)]$ は，収益 $Z^\pi(s,a)=\sum_t\gamma^t R_t$ の期待値のみを扱う．
しかし，金融リスク管理をはじめ多くの意思決定問題では，収益の「平均」よりも「稀に発生する大きな損失」を正確に把握することが重要である．
たとえばオプションを売却した投資家は，原資産を売買してデルタヘッジを行うが，離散的なリヘッジと取引コストのもとでは
ヘッジ損益の分布は非対称で，下側に厚い裾をもつ（正規分布より裾が重く，稀だが巨大な損失が無視できない確率で発生する分布，あるいは損失側に局所的な確率の偏り（多峰性）をもつ歪んだ分布となる）．
このとき投資家が重視すべきは分散ではなく，最悪側の累積確率 $\alpha$ における平均損失である．本レポートでは以下，損益 $Z$ を基準としてこの量を「下側 CVaR$_\alpha$ の\textbf{最大化}」として統一的に扱う（第\ref{sec:drm}節の符号の扱いを参照）\cite{rockafellar2000cvar}．

分布型強化学習 \cite{bellemare2017c51,dabney2018qrdqn,dabney2018iqn,yang2019fqf} は収益分布 $Z^\pi$ 全体を学習するため，
学習した分布から $\CVaR_\alpha$ を計算してリスク回避的な方策を導くことが原理的に可能である．
実際，IQN \cite{dabney2018iqn} では行動選択時の分位点サンプリングを歪めることで CVaR などのリスク尺度に基づく行動選択を行い，
Moghimi and Ku \cite{moghimi2025beyond} は静的スペクトルリスク尺度を最適化する分布型強化学習を提案している．
しかし，これらの手法において「分布をどの点で近似するか」は，リスク尺度とは無関係に，
分布全体を一様に近似するよう定められている．
$\CVaR_{0.05}$ を評価するために必要なのは分位点関数 $\Finv(\tau)$ の $\tau\in[0,0.05]$ における形状であるが，
$N=16$ 個の一様quantile fractionのうち，内部境界 $\tau_i=i/16$（$i=1,\dots,N-1$）でこの区間に入るものは1つもなく，CVaR 領域 $[0, 0.05]$ は実質的に第1区間の中点原子 $\hat\tau_0=1/32=0.03125$ の1点のみで代表される．
これが本研究の出発点である．\textbf{収益分布の下側裾（損失側の裾）を正確に評価したいなら，
近似の「解像度」を下側裾に集中させるべきである}．

\subsection{本研究の位置づけと貢献}
FQF \cite{yang2019fqf} はquantile fraction（累積確率軸 $[0, 1]$ の区切り位置）$\tau_i$ 自体を状態依存的に学習する代表的な分布型強化学習手法であり，
その配置を一様重み付き1-Wasserstein距離 $\Wone$ の最小化に基づいて決定する．
本研究はこの FQF の Fraction Proposal Network の目的関数を，歪み関数 $g$ で重み付けした $\Wone(\dist Z,\dist\hat{Z})$ に置き換える．
本研究は，CVaR最適方策を学習する新しい強化学習アルゴリズムの提案というよりは，
「任意の歪みリスク尺度に対し，限られた離散原子をどこに配置すればリスク尺度推定に有利か」という表現設計・目的指向量子化の研究として位置づけられる．
現時点での貢献は以下の通りである．
\begin{enumerate}
  \item \textbf{理論・アルゴリズム}：各区間の最適原子が $g$-中央値で与えられること（命題\ref{prop:gmedian}），および歪み重み付き $\Wone$ に対する境界勾配が FQF の境界勾配式の $g'(\tau_i)$ 倍となること（命題\ref{prop:grad}）を導出し，これを分位点値ネットワークに plug-in した擬似勾配更新（stop-gradient 型）として実装可能な閉形式を与えた（第\ref{sec:proposed}節）．あわせて，既知の Lipschitz 評価の重み付き形として $|\rho_g(Z)-\rho_g(\hat Z)|\le \Wone(\dist Z,\dist \hat Z)$ が成り立つことを示した（第\ref{sec:drm}節）．
  \item \textbf{最適量子化としての解釈}：FPN の目的関数の最適化が，歪み分布 $\dist Z$ の1次元 $L^1$ 最適量子化問題 \cite{graf2000quantization} と数学的に同値であることを明らかにし，quantile fraction の再配置を目的指向量子化として位置づけた（第\ref{sec:proposed}節）．
  \item \textbf{実験A}：真の分位点関数が解析的に与えられた1次元設定で，$\alpha N\ge2$ 領域における最小改善比 $R=42.3$（$N=16, \alpha=0.05$ では $R=39.6$〜$97.9$）を達成し，原理的有効性を確認した（第\ref{sec:exp}節）．
  \item \textbf{実験B}：裾の重い損失を持つ小規模MDP環境上で実装し，$g$-一様初期化を施した提案法（B1g）において CVaR 推定 MAE が FQF（B0）の約7分の1（$7.02\to0.96$）に低減することを確認した（第\ref{sec:exp}節）．ただしリスク感応 IQN がさらに小さい誤差（$0.36$）を示し，提案法の一般的優位性は未確認である．さらに固定配置の ablation（B3g）により，小規模環境での改善の主因がリスク尺度に整合した配置と原子選択にあることを示し，FPN による状態依存化の追加的有効性は状態依存環境で検証する課題とした．また，greedy 方策の違いによる交絡を共通方策下の交差評価で検討したところ，all-RISKY 分布（重い裾）の参照では全評価器の誤差が大きく，推定対象の難易度が誤差を支配することが示唆された．ただし，本評価は学習時の継続方策の不一致を含むため，表現精度のみの厳密な分離には方策評価モードでの再学習を要する（第\ref{subsec:cross}，第\ref{sec:future}節）．
\end{enumerate}
本レポートは中間報告であり，デルタヘッジ環境での実験と誤差評価の理論的精密化は今後の課題である（第\ref{sec:future}節）．
%==================== 2 分布型強化学習の基礎 ====================
\section{分布型強化学習の基礎}
\label{sec:drl}

\subsection{収益分布と分布ベルマン作用素}
マルコフ決定過程 $(\mathcal{S},\mathcal{A},P,R,\gamma)$ と方策 $\pi$ に対し，収益を確率変数
\begin{equation}
  Z^\pi(s,a)=\sum_{t=0}^{\infty}\gamma^t R_t,\qquad s_0=s,\ a_0=a,\ a_t\sim\pi(\cdot|s_t)
\end{equation}
と定める．$Q^\pi(s,a)=\E[Z^\pi(s,a)]$ である．Bellemare ら \cite{bellemare2017c51} は分布ベルマン作用素
\begin{equation}
  (\mathcal{T}^\pi Z)(s,a)\overset{D}{=}R(s,a)+\gamma Z(S',A'),\qquad S'\sim P(\cdot|s,a),\ A'\sim\pi(\cdot|S')
\end{equation}
を導入し（$\overset{D}{=}$ は分布として等しいことを表す），これが最大形式の $p$-Wasserstein 距離 $\bar d_p(Z_1,Z_2)=\sup_{s,a}W_p(Z_1(s,a),Z_2(s,a))$（$\sup$ は全状態・行動対に対する上限）に関して
$\gamma$-縮小写像であることを示した（なおこの縮小性は，割引率 $\gamma<1$ の下で意思決定が無限ステップ続く無限期間設定（無限ホライズン）で成立する理論であり，第\ref{sec:exp}節の実験Bで扱う最大ステップ数を固定した有限期間設定（有限ホライズン，$\gamma=1$）とは問題構造が異なる点に留意されたい）．
ここで1次元分布 $Z_1,Z_2$ の間の $p$-Wasserstein距離は分位点関数（累積分布関数の一般化逆関数）$F^{-1}_{Z}(\tau)=\inf\{z:F_Z(z)\ge\tau\}$ を用いて
\begin{equation}
  W_p(Z_1,Z_2)=\left(\int_0^1\left|F_{Z_1}^{-1}(\tau)-F_{Z_2}^{-1}(\tau)\right|^p d\tau\right)^{1/p}
  \label{eq:wp}
\end{equation}
と表される．本研究で中心となるのは $p=1$ の場合
\begin{equation}
  \Wone(Z,\hat Z)=\int_0^1\left|\Finv(\tau)-\Fhinv(\tau)\right|d\tau
  \label{eq:w1}
\end{equation}
であり，これを以後「一様重み $\Wone$」と呼ぶ．積分測度が $[0,1]$ 上の一様測度 $d\tau$ であることを強調するためである．

\subsection{C51：固定サポート上のカテゴリカル分布}
C51 \cite{bellemare2017c51} は収益のサポート（分布が確率質量をもつ値の範囲）$[V_{\min},V_{\max}]$ 上に，等間隔の $N=51$ 個の離散的な点
$z_1<\dots<z_N$ を固定し，各点の確率 $p_i(s,a)$ をソフトマックスで出力する．
このように確率質量を置く離散的な点を\textbf{原子（アトム，atom）}と呼ぶ．
ベルマン更新後の分布 $R+\gamma Z'$ の原子は元のサポート上に乗らないため，隣接原子に確率質量を線形分配する射影 $\Phi$ を施し，
射影後の分布との交差エントロピーを最小化する．
理論上は Wasserstein 距離での縮小性が示されているのに対し，実際のアルゴリズムでは KL ダイバージェンスを損失関数に用いるという，理論と実装の乖離が存在した．
またサポート範囲を事前に設定する必要があり，その範囲外の収益が生じる場合には端点へのクリッピング誤差が発生する．
このため，裾の重い収益分布ではサポート範囲の設定に注意を要する．

\subsection{QR-DQN：分位点回帰による固定quantile fraction表現}
QR-DQN \cite{dabney2018qrdqn} は原子の位置ではなく確率側を固定する．
すなわち，各原子に均等な確率 $1/N$ を割り当て，ネットワークがその位置 $\theta_i(s,a)$ を出力することで，$Z(s,a)\approx\frac1N\sum_{i=1}^N\delta_{\theta_i(s,a)}$ と近似する．ここで $\delta_\theta$ は点 $\theta$ に確率質量1を置く離散分布（ディラック測度）である．
Dabney ら \cite{dabney2018qrdqn} は，等確率 $1/N$ の $N$ 点分布のうち $\Wone(Z,\hat Z)$ を最小化する離散分布の原子が，分位点中点 $\hat\tau_i=(2i-1)/(2N)$ における分位点値
$\theta_i=\Finv(\hat\tau_i)$ で与えられることを示した．QR-DQN はこの原子位置を分位点回帰で学習する．
分位点 $\tau$ に対する分位点回帰損失は，指示関数 $\mathbb{I}$ を用いて
\begin{equation}
  \rho_\tau(u)=u\left(\tau-\mathbb{I}\{u<0\}\right),\qquad u=\text{（目標）}-\text{（予測）}
\end{equation}
であり，その期待値の最小点が $\tau$-分位点になる．
誤差ゼロ付近での非微分性（滑らかでないこと）に起因する学習の不安定性を避けるため，実装では Huber 損失 $\mathcal{L}_\kappa$ を組み合わせた
\begin{equation}
  \rho^\kappa_\tau(u)=\left|\tau-\mathbb{I}\{u<0\}\right|\frac{\mathcal{L}_\kappa(u)}{\kappa},\qquad
  \mathcal{L}_\kappa(u)=\begin{cases}\tfrac12u^2 & |u|\le\kappa\\ \kappa(|u|-\tfrac12\kappa)& \text{それ以外}\end{cases}
  \label{eq:huber-quantile}
\end{equation}
を用いる．目標分布の原子 $\theta'_j$ と予測原子 $\theta_i$ のすべての対について
\begin{equation}
  \mathcal{L}_{\mathrm{QR}}=\frac1N\sum_{i=1}^N\sum_{j=1}^N\rho^\kappa_{\hat\tau_i}\!\left(r+\gamma\theta'_j-\theta_i\right)
  \label{eq:qr-loss}
\end{equation}
を最小化する．この定式化により損失が Wasserstein 距離の最小化と整合的になり，サポートの事前固定も不要になった．
なお，式\eqref{eq:qr-loss}は QR-DQN や IQN における標準形（等重み $1/N$ による平均）を示しているが，第\ref{sec:proposed}節で述べる FQF 系ではターゲット原子が区間幅に応じて重み付けされる点に留意されたい．
本研究の観点から重要なのは，QR-DQN の近似能力が「quantile fraction が等間隔に固定されていること」によって制約される点である．

\subsection{IQN：暗黙的分位点ネットワーク}
\label{subsec:iqn}

\subsubsection{QR-DQNの限界と動機：離散表現の制約}
前項で述べた QR-DQN \cite{dabney2018qrdqn} は，原子の確率を $1/N$ に固定してその位置を分位点回帰で推定することで，C51 のサポート事前固定問題を克服した．
しかし，強化学習の表現設計の観点からは，依然として固定された格子による離散性と計算コストに関する制約が残されていた．
% (重複行を前文に統合)

QR-DQN は，あらかじめ決められた $N$ 個（例えば $N=32$ や $N=200$）の分位点中点 $\hat\tau_i=(2i-1)/(2N)$ における値しか出力できない．
例えば意思決定において累積確率 $\tau=0.03$ 等の任意の分位点値を評価する場合，QR-DQN は任意の $\tau$ に対する値を直接出力できず，隣接する格子点から線形補間する必要がある．
さらに，解像度を高めようとして原子数 $N$ を増やすと，ネットワークの最終出力層の次元数 $N \times |\mathcal{A}|$ が肥大化し，パラメータ数と計算コストが増大する．

分位点関数 $\Finv(\tau)$ は本来，累積確率軸 $[0, 1]$ 上で定義された連続な1次元関数である．
したがって，離散的な $N$ 個の値を出力するのではなく，\textbf{状態 $s$ に加えて任意の分位点レベル $\tau \in [0, 1]$ もネットワークの入力とし，その $\tau$ に対する分位点値 $\Finv(\tau)$ を直接出力する}という連続関数としてのモデリングが着想される．
この着想に基づき，Dabney ら \cite{dabney2018iqn} は \textbf{暗黙的分位点ネットワーク（Implicit Quantile Networks; IQN）} を提案した．
「暗黙的（implicit）」と呼ばれるのは，分布のパラメータ（平均・分散や固定原子の確率など）を明示的に出力するのではなく，任意の $\tau$ を入力して対応する値を取り出す「入出力関係（生成写像）」として分布を表現しているためである．

\subsubsection{IQNのアーキテクチャとコサイン基底埋め込み}
IQN は，状態 $s$ と分位点レベル $\tau \sim U(0, 1)$ を入力とし，行動 $a$ の分位点値 $Z_\tau(s, a)$ を出力する．
ここで，1次元の実数値である $\tau \in [0, 1]$ を状態特徴ベクトル $\psi(s)$ に直接結合せず，
コサイン基底の級数に展開して埋め込む理由について述べる．

スカラ値 $\tau$ を直接入力する構成が不十分である背景には，
深層ニューラルネットワークの性質である \textbf{周波数バイアス（Spectral Bias）} \cite{rahaman2019spectral} がある．
標準的な MLP は，入力に対して低周波の滑らかな関数を優先的に学習し，高周波成分や急激な変化の学習が遅いという性質をもつ．
しかし，確率分布の分位点関数 $\Finv(\tau)$ は，裾（$\tau \to 0$ や $\tau \to 1$）において発散傾向を示し，あるいは局所的な確率の偏りにおいて階段状の急な立ち上がりを示す．
スカラ値 $\tau$ をそのまま入力したのでは，MLP はその強い非線形性を捉えきれず，全体として緩やかな関数を出力してしまい，分布の形状が平滑化されてしまう．

そこで IQN では，\textbf{コサイン基底を用いた高次元埋め込み（コサイン埋め込み）} を導入する．
具体的には，$n$ 次のコサイン基底 $\{\cos(\pi k \tau)\}_{k=0}^{n-1}$ を用い，埋め込みベクトルの第 $j$ 成分を
\begin{equation}
  \phi_j(\tau) = \mathrm{ReLU}\!\left(\sum_{k=0}^{n-1}\cos(\pi k\tau)w_{jk} + b_j\right),\qquad j=1,\dots,d
  \label{eq:cosine-emb}
\end{equation}
と定義する（ここで $w_{jk}, b_j$ は学習可能な重みとバイアス，$d$ は埋め込み次元である）．
$k=0$ は定数項であり，$k=1, 2, \dots$ と周波数が上がるにつれて，入力 $\tau$ の微小な変化に反応する高周波基底となる．
この高次元埋め込みベクトル $\phi(\tau) \in \R^d$ を，状態 $s$ の特徴量ベクトル $\psi(s) \in \R^d$ と \textbf{要素積（Hadamard積 $\psi(s) \odot \phi(\tau)$）} によって結合する．
要素積を採用する理由は，状態特徴 $\psi(s)$ の各次元が各周波数のコサイン基底の振幅を変調する相互作用を表現できるためである．
こうして状態と分位点レベルが融合された表現が後続の全結合層を通り，最終的な分位点値 $Z_\tau(s, a)$ が計算される．

\subsubsection{IQNの学習損失}
IQN の学習では，固定された格子点は存在しない．
リプレイバッファから取り出した各遷移 $(s, a, r, s')$ に対し，現在状態の予測用として $N$ 個の分位点レベル $\tau_i \sim U(0, 1)$（$i=1,\dots,N$）を独立にサンプリングし，次状態のターゲット作成用として $N'$ 個の分位点レベル $\tau'_j \sim U(0, 1)$（$j=1,\dots,N'$）を独立にサンプリングする．
% (前行に統合)
% (前行に統合)
そして，QR-DQN と同型の Huber 分位点回帰損失（式\eqref{eq:huber-quantile}）を用いて，すべての対 $(i, j)$ にわたる標本平均
\begin{equation}
  \mathcal{L}_{\mathrm{IQN}} = \frac{1}{N N'}\sum_{i=1}^N\sum_{j=1}^{N'}\rho^\kappa_{\tau_i}\!\left(r + \gamma Z_{\tau'_j}(s', a^*) - Z_{\tau_i}(s, a)\right)
  \label{eq:iqn-loss}
\end{equation}
を最小化する（ここで $a^* = \arg\max_{a'} \frac{1}{K}\sum_{k=1}^K Z_{\tau_k}(s', a')$ は次状態における貪欲行動である）．
更新のたびに異なる乱数 $\tau$ がサンプリングされるため，ネットワークは訓練が進むにつれて $[0, 1]$ の全域において連続的かつ滑らかな分位点関数 $\Finv(\tau)$ を学習していく．

\subsubsection{行動選択におけるリスク感応性と，その構造的課題}
IQN の大きな利点は，\textbf{学習後の行動選択において $\tau$ のサンプリング分布を変更するだけで，追加の学習を行うことなく任意の歪みリスク尺度に基づく行動決定を実行できる} 点にある．
行動選択時の $\tau$ のサンプリング分布を $\beta$（リスクプロファイル）と表すと，方策は
\begin{equation}
  \pi_\beta(s) = \arg\max_a\ \E_{\tau\sim\beta}\left[Z_\tau(s,a)\right]
  \label{eq:iqn-risk-policy}
\end{equation}
と定義される．
とくに，累積確率の下側 $\alpha$（下側裾）におけるリスクを重視する場合，サンプリング分布として区間 $[0, \alpha]$ 上の一様分布 $\beta = U(0, \alpha)$ を設定する．
このとき，
\begin{equation}
  \E_{\tau\sim U(0,\alpha)}\left[Z_\tau(s, a)\right] = \frac{1}{\alpha}\int_0^\alpha \Finv(u; s, a)\,du = \CVaR_\alpha(Z(s, a))
\end{equation}
となり，モンテカルロ平均をとるだけで，定義通りの条件付きバリューアットリスク（$\CVaR_\alpha$）に基づいたリスク回避的行動選択が実現できる．
これは第\ref{sec:drm}節で後述する歪みリスク尺度の考え方と一致している．

しかし，この定式化には表現リソースの配分に関する課題が存在する．
IQN の標準的な定式化では，
\begin{itemize}
  \item \textbf{学習時}：$\tau \sim U(0, 1)$ から一様にサンプリングし，分布全体を均等に学習する．
  \item \textbf{行動選択時}：下側裾 $\tau \sim U(0, \alpha)$ だけを切り取って意思決定を行う．
\end{itemize}
このように，学習時と行動選択時でサンプリング方針が分離している．

もし最初から「最悪損失の低減（下側 $\CVaR_\alpha$ の最大化）」が目的であると分かっているならば，この枠組みには重大な非効率性が潜んでいる．
仮に $\alpha=0.05$ だとすると，エージェントの意思決定に使われるのは累積確率軸のわずか $5\%$（$[0, 0.05]$）にすぎない．
にもかかわらず，ネットワークの限られたパラメータと表現容量の大半は，意思決定において全く使われない残り $95\%$（$[0.05, 1]$）の領域を一様に近似するために費やされているのである．
下側裾のリスク評価を目的とする場合，学習段階から下側裾に表現力を集中させることが望ましいという問題意識が生じる．
この課題意識に対する解決策が FQF の登場であり，本研究の提案手法へとつながる．

%--------------------------------------------------
\subsection{FQF：quantile fractionの完全パラメータ化}
\label{subsec:fqf}

\subsubsection{FQFの着眼点：なぜquantile fractionを学習するのか}
IQN は連続的な分位点関数を扱える利点をもつ一方，意思決定やベルマンターゲットの計算のたびに $\tau$ をランダムサンプリングするため，推定値にサンプリング分散（ノイズ）が伴う．
一方，QR-DQN は決定論的であるものの，等間隔に固定された格子点に制約されていた．

もし，固定された等間隔格子でもなく，ランダムサンプリングでもなく，
状態に応じて累積確率軸 $[0, 1]$ の区切り位置そのものを適応的に学習できれば，表現力をより効率的に活用できると考えられる．
値の変化が大きい重要な領域（下側裾など）には区切りを密に配置して解像度を高め，変化の緩やかな領域には粗く区切りを配置することで，
限られた $N$ 個の代表点で分布の近似精度を高める手法として，Yang ら \cite{yang2019fqf} による \textbf{FQF（Fully Parameterized Quantile Function）} が提案された．

\subsubsection{FPNと分位点値ネットワークの二段構造}
FQF は，累積確率軸 $[0, 1]$ を $N$ 個の区間に分割する境界点 $0 = \tau_0 < \tau_1 < \dots < \tau_{N-1} < \tau_N = 1$ を状態依存のパラメータとして学習する．
この区切り位置 $\tau_i$ を本レポートでは一貫して \textbf{quantile fraction} と呼ぶ．
ネットワークは役割の異なる2つの構成要素からなる．

\begin{enumerate}
  \item \textbf{Fraction Proposal Network（FPN）} $\omega_1$：
  状態特徴 $\psi(s)$ を入力とし，$N$ 個の実数を出力する．これをソフトマックス関数に通すことで，各区間の幅を表す確率ベクトル $p = (p_0, p_1, \dots, p_{N-1}) \in \Delta^{N-1}$（$\Delta^{N-1}$ は各成分が正で総和が1となる単体）を得る．
  そして，この区間幅の累積和（cumulative sum）をとることによって，
  \begin{equation}
    \tau_0 = 0,\qquad \tau_i = \sum_{k=0}^{i-1} p_k\quad (i=1,\dots,N-1),\qquad \tau_N = 1
    \label{eq:fpn-cumsum}
  \end{equation}
  として quantile fraction 列を生成する．
ソフトマックスと累積和を組み合わせる設計により，特別な制約処理を要さずに，\textbf{常に狭義単調増加 $0 = \tau_0 < \tau_1 < \dots < \tau_N = 1$ が保証される}．

  \item \textbf{Quantile Value Network} $\omega_2$：
  IQN と同様にコサイン基底埋め込みを備えたネットワークであり，各区間 $[\tau_i, \tau_{i+1}]$ を代表する分位点レベル $\hat\tau_i$ を入力として，その点における分位点値 $\Finv_{\omega_2}(\hat\tau_i; s, a)$ を出力する．
\end{enumerate}

こうして得られる近似分布は，各点 $\Finv(\hat\tau_i)$ に区間幅 $\tau_{i+1}-\tau_i$ の確率質量を置いた離散分布
\begin{equation}
  \hat Z_N = \sum_{i=0}^{N-1}(\tau_{i+1}-\tau_i)\,\delta_{\Finv(\hat\tau_i)}
  \label{eq:zhat_detailed}
\end{equation}
である（ここで $\delta_z$ は点 $z$ に確率1をもつディラック測度）．
対応する近似分位点関数 $\Fhinv(u)$ は，各区間 $u \in (\tau_i, \tau_{i+1}]$ 上で一定値 $\Finv(\hat\tau_i)$ をとる階段関数となる．
\paragraph{FPNの動作の数値具体例}
理解を具体化するため，区間数 $N=4$ の数値例を考える．
状態 $s$ において FPN の出力層が $(2.0, 1.0, 0.0, 0.0)$ という実数ベクトルを出力したとする．
これをソフトマックス関数に通すと，各成分が正で総和が1となる区間幅ベクトル
$p \approx (0.505, 0.186, 0.068, 0.241)$
が得られる（値の変化が激しい第1区間に最も広い確率幅が割り振られている）．
この累積和をとることで，境界列
$\tau_0 = 0,\quad \tau_1 = 0.505,\quad \tau_2 = 0.691,\quad \tau_3 = 0.759,\quad \tau_4 = 1.0$
が厳密な狭義単調増加列として自動的に定まる．
各区間の算術中点は
$\hat\tau_0 \approx 0.253,\quad \hat\tau_1 \approx 0.598,\quad \hat\tau_2 \approx 0.725,\quad \hat\tau_3 \approx 0.880$
となり，分位点値ネットワーク $\omega_2$ はこの4点における値を評価する．
このように，ソフトマックスと累積和の組み合わせによって単調増加性が保証される．


\subsubsection{補題1（一様重み $W_1$ に対する最適原子）の証明}
ここで，境界 $\tau_0 < \tau_1 < \dots < \tau_N$ が与えられたとき，各区間 $[\tau_i, \tau_{i+1}]$ を代表する原子の位置 $q_i$ の最適配置に関する問題が生じる．
% 各区間の代表原子配置の最適性

Yang ら \cite{yang2019fqf} の Lemma 1 は，真の分位点関数 $\Finv$ に対する一様重み 1-Wasserstein 距離 $\Wone(Z, \hat Z_N) = \int_0^1 |\Finv(u) - \Fhinv(u)| du$ を最小化する最適な原子の位置が，\textbf{各区間の算術中点 $\hat\tau_i = (\tau_i + \tau_{i+1})/2$ における分位点値 $\Finv(\hat\tau_i)$} で与えられることを示した．
以下，この証明の各ステップを確認する．

\begin{lemma}[一様重み $W_1$ に対する最適原子]
\label{lem:fqf-atom}
任意の広義単調増加な連続分位点関数 $\Finv$ と，固定された quantile fraction $\tau_i < \tau_{i+1}$ に対し，
区間損失
\begin{equation}
  \ell_i(q) = \int_{\tau_i}^{\tau_{i+1}} |\Finv(u) - q|\,du
  \label{eq:lem1-loss}
\end{equation}
を最小化する原子 $q \in \R$ は，中点 $\hat\tau_i = (\tau_i + \tau_{i+1})/2$ の分位点値 $q^* = \Finv(\hat\tau_i)$ である．
\end{lemma}

\begin{proof}
目的関数 $\ell_i(q)$ は絶対値を含む積分であるため，まず被積分関数の絶対値を展開する．
$\Finv$ は連続かつ非減少であるから，$q = \Finv(m)$ を満たす点 $m \in [\tau_i, \tau_{i+1}]$ が存在する（もし $q < \Finv(\tau_i)$ や $q > \Finv(\tau_{i+1})$ ならば明らかに損失が増大するため，最適解は区間の内部に存在する）．
この点 $m$ によって積分区間を分割すると，
\begin{itemize}
  \item $u \in [\tau_i, m]$ においては，$\Finv$ の単調性より $\Finv(u) \le \Finv(m) = q$ であるから，$|\Finv(u) - q| = q - \Finv(u)$ となる．
  \item $u \in [m, \tau_{i+1}]$ においては，同様に $\Finv(u) \ge \Finv(m) = q$ であるから，$|\Finv(u) - q| = \Finv(u) - q$ となる．
\end{itemize}
したがって，区間損失 $\ell_i(q)$ は絶対値記号を用いずに次のように表される．
\begin{equation}
  \ell_i(q) = \int_{\tau_i}^{m(q)} \left(q - \Finv(u)\right)du + \int_{m(q)}^{\tau_{i+1}} \left(\Finv(u) - q\right)du.
  \label{eq:lem1-split}
\end{equation}
この式を $q$ で微分して極値を求める．
ここで，積分の上限・下限にある $m(q)$ もまた $q$ に依存して変化する（$q = \Finv(m)$ より $dm/dq = 1/(\Finv)'(m)$）ことに留意されたい．
一般に，積分記号下の微分に関する \textbf{Leibniz（ライプニッツ）の積分則}は次式で与えられる．
\begin{equation}
  \frac{d}{dq}\left[\int_{a(q)}^{b(q)} f(u, q)\,du\right] = f(b(q), q)\,\frac{db}{dq} - f(a(q), q)\,\frac{da}{dq} + \int_{a(q)}^{b(q)} \frac{\partial f}{\partial q}(u, q)\,du
  \label{eq:leibniz-general}
\end{equation}
を式\eqref{eq:lem1-split}の各項に適用する．すると，
\begin{align}
  \frac{d}{dq}\left[\int_{\tau_i}^{m(q)} \left(q - \Finv(u)\right)du\right]
  &= \underbrace{\left(q - \Finv(m)\right)}_{=\,0}\,\frac{dm}{dq} - 0 + \int_{\tau_i}^m \frac{\partial}{\partial q}(q - \Finv(u))\,du \notag\\
  &= \int_{\tau_i}^m 1\,du = m - \tau_i,
  \label{eq:leibniz-term1}\\
  \frac{d}{dq}\left[\int_{m(q)}^{\tau_{i+1}} \left(\Finv(u) - q\right)du\right]
  &= 0 - \underbrace{\left(\Finv(m) - q\right)}_{=\,0}\,\frac{dm}{dq} + \int_m^{\tau_{i+1}} \frac{\partial}{\partial q}(\Finv(u) - q)\,du \notag\\
  &= \int_m^{\tau_{i+1}} (-1)\,du = -( \tau_{i+1} - m ).
  \label{eq:leibniz-term2}
\end{align}
ここで重要なのは，\textbf{境界点 $u=m$ において被積分関数 $(q - \Finv(u))$ および $(\Finv(u) - q)$ がゼロとなるため，端点 $m(q)$ の変動に伴う微分項が厳密にゼロとなる} という点である．
したがって，足し合わせた導関数は次のように整理される．
\begin{equation}
  \frac{d\ell_i}{dq} = (m - \tau_i) - (\tau_{i+1} - m) = 2m - (\tau_i + \tau_{i+1}).
  \label{eq:lem1-deriv}
\end{equation}
損失が最小となる最適条件（停留点条件）は $\frac{d\ell_i}{dq} = 0$ であるから，
\begin{equation}
  2m - (\tau_i + \tau_{i+1}) = 0 \iff m = \frac{\tau_i + \tau_{i+1}}{2} = \hat\tau_i.
\end{equation}
さらに2階微分を計算すると $\frac{d^2\ell_i}{dq^2} = 2 \frac{dm}{dq} = \frac{2}{(\Finv)'(m)} > 0$（$\Finv$ が狭義単調増加のとき）となり，凸関数の最小点であることが確認される．
よって，区間損失を最小化する原子は算術中点の分位点値 $q^* = \Finv(\hat\tau_i)$ に一意に定まる．
\end{proof}

この導出結果から，直観的な幾何学的解釈が得られる．
式\eqref{eq:lem1-deriv}の導関数 $(m - \tau_i) - (\tau_{i+1} - m)$ は，点 $m$ の左側にある区間の長さと，右側にある区間の長さの「差」を表している．
導関数がゼロになる点とは，左側の長さと右側の長さが釣り合う点，すなわち \textbf{区間の一様測度に関する中央値（median）} にほかならない．
一様測度においては中央値が算術中点と一致するため $\hat\tau_i = (\tau_i+\tau_{i+1})/2$ が得られたのである．
この「測度の中央値」という観点は，第\ref{sec:proposed}節で述べる歪みリスク尺度への拡張の基礎となる．

\subsubsection{命題1（quantile fractionに関する勾配の閉形式）の導出と包絡線定理}
各区間における最適原子の位置が中点 $q_i = \Finv(\hat\tau_i)$ として定まったところで，
次に FPN のパラメータを更新するための境界勾配を導出する．
全体の近似損失は，全区間の損失の総和
\begin{equation}
  \Wone(Z, \hat Z_N) = \sum_{j=0}^{N-1} \ell_j(q_j) = \sum_{j=0}^{N-1} \int_{\tau_j}^{\tau_{j+1}} |\Finv(u) - q_j|\,du
  \label{eq:total-w1}
\end{equation}
である．これを内部の境界 $\tau_i$（$i=1,\dots,N-1$）で偏微分したい．
Yang ら \cite{yang2019fqf} の Proposition 1 は，この勾配が次のように簡潔な閉形式で表現できることを示している．

\begin{proposition}[一様 $\Wone$ に対する境界勾配]
\label{prop:fqf-grad}
任意の $i=1,\dots,N-1$ に対し，
\begin{equation}
  \frac{\partial \Wone(Z, \hat Z_N)}{\partial\tau_i} = 2\Finv(\tau_i) - \Finv(\hat\tau_i) - \Finv(\hat\tau_{i-1}).
  \label{eq:fqf-prop1}
\end{equation}
\end{proposition}

\begin{proof}
総和\eqref{eq:total-w1}の中で，境界 $\tau_i$ に依存する項は区間 $i-1$（$[\tau_{i-1}, \tau_i]$）の積分上限および区間 $i$（$[\tau_i, \tau_{i+1}]$）の積分下限のみである．
他の区間 $j \notin \{i-1, i\}$ は $\tau_i$ を含まないため微分はゼロとなる．
% (重複削除)
したがって，
\begin{equation}
  \frac{\partial \Wone}{\partial\tau_i} = \frac{\partial}{\partial\tau_i}\left[ \ell_{i-1}(q_{i-1}) + \ell_i(q_i) \right]
\end{equation}
を計算すればよい．

合成関数の偏微分律に従うと，$\tau_i$ に関して2つの経路が存在する．
\begin{itemize}
  \item \textbf{経路(a)}：積分区間の端点（境界）の変動に伴う項．
  \item \textbf{経路(b)}：最適原子 $q_{i-1}, q_i$ の配置の変動に伴う項．
\end{itemize}
すなわち，全微分は
\begin{equation}
  \frac{\partial \Wone}{\partial\tau_i} = \underbrace{\frac{\partial \ell_{i-1}}{\partial\tau_i}\Bigg|_{q_{i-1}\text{固定}} + \frac{\partial \ell_i}{\partial\tau_i}\Bigg|_{q_i\text{固定}}}_{\text{経路(a): 端点の移動による寄与}} + \underbrace{\frac{\partial \ell_{i-1}}{\partial q_{i-1}}\,\frac{\partial q_{i-1}}{\partial\tau_i} + \frac{\partial \ell_i}{\partial q_i}\,\frac{\partial q_i}{\partial\tau_i}}_{\text{経路(b): 原子の移動による寄与}}
  \label{eq:chain-rule-grad}
\end{equation}
となる．

ここで，最適化理論における\textbf{包絡線定理（Envelope Theorem）}を適用する．
ここで，補題\ref{lem:fqf-atom}により，
原子 $q_{i-1}$ および $q_i$ は，それぞれ区間損失 $\ell_{i-1}(q)$ および $\ell_i(q)$ を最小化するように選ばれている．
最小点における停留条件により，
\begin{equation}
  \frac{\partial \ell_{i-1}}{\partial q_{i-1}} = 0,\qquad \frac{\partial \ell_i}{\partial q_i} = 0
\end{equation}
が厳密に成立している．
したがって，$\frac{\partial q_{i-1}}{\partial\tau_i} \ne 0$ や $\frac{\partial q_i}{\partial\tau_i} \ne 0$ であっても，それにゼロが乗じられるため，\textbf{経路(b)の項は厳密にゼロとなり消滅する．}
「最適化された内部パラメータの変動による目的関数の変化率は，最適解自身の変動を考慮せず直接の偏微分のみで評価できる」という包絡線定理により，最適原子の配置変動に伴う微分を考慮する必要がなくなる．

残る経路(a)の端点微分を計算する．
Leibnizの積分則により被積分関数に端点 $\tau_i$ を代入すると，区間 $i-1$ の積分上限に関する項から $+|\Finv(\tau_i) - q_{i-1}|$，区間 $i$ の積分下限に関する項から $-|\Finv(\tau_i) - q_i|$ が得られる．
% (前文に統合)
% (前文に統合)
したがって，
\begin{equation}
  \frac{\partial \Wone}{\partial\tau_i} = |\Finv(\tau_i) - q_{i-1}| - |\Finv(\tau_i) - q_i|
  \label{eq:abs-grad}
\end{equation}
である．
ここで，中点 $\hat\tau_{i-1}$ と境界 $\tau_i$ と中点 $\hat\tau_i$ の大小関係は明らかに $\hat\tau_{i-1} < \tau_i < \hat\tau_i$ である．
$\Finv$ は広義単調増加であるから，
\begin{equation}
  q_{i-1} = \Finv(\hat\tau_{i-1}) \le \Finv(\tau_i) \le \Finv(\hat\tau_i) = q_i
\end{equation}
が成り立つ．よって絶対値記号はそのまま外すことができ，
\begin{align}
  |\Finv(\tau_i) - q_{i-1}| &= \Finv(\tau_i) - \Finv(\hat\tau_{i-1}), \\
  |\Finv(\tau_i) - q_i| &= \Finv(\hat\tau_i) - \Finv(\tau_i)
\end{align}
となる．これを式\eqref{eq:abs-grad}に代入すると，
\begin{align}
  \frac{\partial \Wone}{\partial\tau_i}
  &= \left(\Finv(\tau_i) - \Finv(\hat\tau_{i-1})\right) - \left(\Finv(\hat\tau_i) - \Finv(\tau_i)\right) \notag\\
  &= 2\Finv(\tau_i) - \Finv(\hat\tau_i) - \Finv(\hat\tau_{i-1})
\end{align}
以上により，一様境界勾配式\eqref{eq:fqf-prop1}が導出される．
\end{proof}

\subsubsection{勾配式の直観的解釈と最適量子化（Lloyd の条件）}
式\eqref{eq:fqf-prop1}は次のような幾何学的意味をもつ．
勾配降下法において，境界 $\tau_i$ の更新が停止する停留点条件 $\frac{\partial \Wone}{\partial\tau_i} = 0$ は，
\begin{equation}
  2\Finv(\tau_i) - \Finv(\hat\tau_i) - \Finv(\hat\tau_{i-1}) = 0 \iff \Finv(\tau_i) = \frac{\Finv(\hat\tau_{i-1}) + \Finv(\hat\tau_i)}{2}
  \label{eq:lloyd-boundary}
\end{equation}
と書ける．
この等式は，
\textbf{「境界 $\tau_i$ における分位点値 $\Finv(\tau_i)$ が，左隣の原子 $\Finv(\hat\tau_{i-1})$ と右隣の原子 $\Finv(\hat\tau_i)$ の算術平均（中点）にあるとき勾配がゼロになる」}
という幾何学的関係を表している．

$\Finv(\tau_i)$ が右隣の原子 $\Finv(\hat\tau_i)$ に偏っている場合，勾配は正となり境界 $\tau_i$ を減少させる方向に働く．
逆に左隣の原子に偏っている場合は勾配が負となり境界を増加させる方向に働く．
この勾配がゼロとなる平衡点が式\eqref{eq:lloyd-boundary}である．

これは，クラスタリングの \textbf{$k$-means 法（Lloyd のアルゴリズム）} と同等の構造である．
\begin{enumerate}
  \item \textbf{境界の条件}：各区間の境界は，隣接する代表点の中点で区切られる（Voronoi 分割）．
  \item \textbf{代表点の条件}：各クラスタの代表点は，そのクラスタに属するデータ点の中央値（$L^1$ ノルムの場合）に置かれる．
\end{enumerate}
FQF の目的関数を一様 $\Wone$ 距離で最適化することは，\textbf{「連続確率分布 $Z$ を，1次元空間において $N$ 個の代表点に最も誤差が小さくなるように離散化（クラスタリング）する問題」}，すなわち数学における \textbf{1次元 $L^1$ 最適量子化問題（Optimal Quantization）} \cite{graf2000quantization} を解くことに相当する．

%--------------------------------------------------
\subsection{系譜のまとめ}
\label{subsec:lineage}

ここまでの分布型強化学習の発展の流れを表\ref{tab:lineage}に整理する．
C51 から FQF への進化は，
\begin{itemize}
  \item \textbf{値側の固定（C51）}：サポート $[V_{\min}, V_{\max}]$ を固定し，確率を学習．
  \item $\to$ \textbf{確率側の固定（QR-DQN）}：確率を等間隔 $1/N$ に固定し，値を分位点回帰で学習．
  \item $\to$ \textbf{確率側の連続化（IQN）}：確率 $\tau \sim U(0, 1)$ を連続的にサンプリングし，任意の分位点を学習．
  \item $\to$ \textbf{確率側の最適化（FQF）}：確率境界 $\tau_i$ をニューラルネットワークで適応的に学習．
\end{itemize}
というように，分布の表現自由度を段階的に拡張してきた系譜として位置づけられる．

\begin{table}[htbp]
\centering
\caption{分布型強化学習の系譜と本研究の位置づけ．最右列が本研究で拡張・置換する部分．}
\label{tab:lineage}
\small
\setlength{\tabcolsep}{2pt}
\begin{tabular}{lllll}
\toprule
手法 & 表現 & 固定するもの & 損失 & quantile fractionの基準\\
\midrule
C51 \cite{bellemare2017c51} & カテゴリカル & サポート原子位置 & 射影後交差エントロピー & ---\\
QR-DQN \cite{dabney2018qrdqn} & 分位点 & quantile fraction $i/N$ & Huber分位点回帰 & 一様固定（$i/N$ 等間隔）\\
IQN \cite{dabney2018iqn} & 暗黙的分位点 & なし（連続入力） & Huber分位点回帰 & 一様サンプリング\\
FQF \cite{yang2019fqf} & 可変分位点 & なし & 分位点回帰＋$\Wone$ & 一様重み $\Wone$ 最小化\\
\midrule
\textbf{提案法（歪みFQF）} & 可変分位点 & なし & 分位点回帰＋$\Wone(\dist Z,\dist\hat Z)$ & \textbf{歪み重み $\Wone$ 最小化}\\
\bottomrule
\end{tabular}
\end{table}

しかし，表\ref{tab:lineage}の最右列に示すように，FQF においても quantile fraction を決定する評価基準は一様重み $\Wone$ 距離（分布全体を一様に近似する基準）にとどまっていた．
金融のリスク管理や安全性が重視される制御のように，大きな損失の回避を目的とする場合，quantile fractionの決定基準自体をリスク尺度に適合させる必要がある．
これが，本研究が提案する \textbf{歪みリスク尺度に基づくquantile fraction学習（歪み FQF）} の動機である．
次章において，その数理的基盤となる歪みリスク尺度と歪み分布について述べる．

%==================== 3 歪みリスク尺度 ====================
\section{歪みリスク尺度}
\label{sec:drm}

\subsection{定義とYaariの双対理論}
以下，$Z$ は損益（profit and loss; 正の値が大きいほど望ましく，負の値が大きいほど損失を意味する実数値確率変数）を表し，リスクは下側裾（マイナス側への極端なブレ）にあるとする．
したがって本研究で「リスク尺度」と呼ぶ $\rho_g(Z)$ は，正の損失 $L=-Z$ を最小化する損失型リスク尺度ではなく，\textbf{収益 $Z$ に対する歪み重み付き評価（risk-adjusted value）} であり，下側 CVaR の\textbf{最大化}を目的とする方針で統一する（金融工学等で損失を正値として定義し最小化する定式化と符号が逆になる点に留意されたい）．

\subsubsection{期待効用理論とYaariの双対理論の違い}
本研究で期待値基準や平均分散基準ではなく歪みリスク尺度を採用する数理的背景について述べる．

通常の強化学習では，収益の期待値 $\E[Z]$ を最大化する．
しかし，期待値基準では「平均利益は正であるが，確率1\%で極めて大きな損失が生じる」ような行動も選択されうる．
現実の意思決定（自動運転，医療，金融リスク管理）では，稀に発生する深刻な損失（極度のヘビーテール）に対して適切に対処しなければならない．

このようなリスク回避的な意思決定を数理モデル化する伝統的な手法として，フォン・ノイマンとモルゲンシュテルンによる \textbf{期待効用理論（Expected Utility Theory）} がある．
期待効用理論では，結果の金額 $z$ を凹関数である効用関数 $u(z)$（例えば $u(z) = 1 - e^{-\gamma z}$）によって主観的な満足度に変換し，その期待値 $\E[u(Z)]$ を評価する．
これは「1億円もらえる嬉しさは，100万円の100倍ではない（限界効用逓減）」という人間の主観的価値を表現している．
しかし，期待効用理論においては，事象の生起確率 $F(z)$ そのものは客観的な確率として忠実に扱われる．

これに対して，Yaari \cite{yaari1987dual} が提唱した \textbf{双対理論（Dual Theory of Choice under Risk）} は，確率の側を歪めるアプローチをとる．
\textbf{「貨幣価値 $z$ そのものは客観的な金額として線形のまま扱い，事象の生起確率（累積確率 $\tau$）の側を人間の心理やリスク選好に応じて非線形に歪める」} のである．
客観的な利得の値ではなく確率の側を歪めるこの枠組みは，数理ファイナンスや保険数理において広く用いられるリスク尺度につながった．それが \textbf{歪みリスク尺度（Distortion Risk Measure; DRM）} である．

\subsubsection{歪みリスク尺度の数学的定義}
\begin{definition}[歪み関数と歪みリスク尺度]
\label{def:drm}
非減少で $g(0)=0$ かつ $g(1)=1$ を満たす連続写像 $g:[0,1]\to[0,1]$ を\textbf{歪み関数}と呼ぶ．
$g$ が絶対連続でルベーグ可測な導関数 $g'$ をもつとき，収益 $Z$ に対する歪みリスク尺度を
\begin{equation}
  \rho_g(Z) = \int_0^1 \Finv(\tau)\,dg(\tau) = \int_0^1 \Finv(\tau)\,g'(\tau)\,d\tau
  \label{eq:drm}
\end{equation}
と定義する．
\end{definition}

なお，本研究で用いる歪み関数は，後述の $g_\alpha$，$g_{\alpha,\lambda}$，$g_\eta$ を含めいずれも区分的になめらかであり，$g'$ は区分定数または連続であるため，上記式の積分は常に適切に定まる．
$g(\tau)=\tau$（恒等写像）のときは $g'(\tau)=1$ となり，$\rho_g(Z) = \int_0^1 \Finv(\tau)d\tau = \E[Z]$ と通常の期待値に一致する．

\paragraph{歪み関数とリスクの直観}
歪み関数 $g$ が\textbf{凹関数}であるとき，その導関数 $g'(\tau)$ は非増加関数となる．
式\eqref{eq:drm}において，$\tau \in [0, 1]$ は累積確率であり，$\tau \to 0$ は最も悪い結果（極端な損失，下側裾）に対応する．
$g'$ が非増加であるということは，\textbf{小さい $\tau$（悪い結果）に対して非常に大きな重み $g'(\tau) \gg 1$ が乗せられ，逆に大きい $\tau$（良い結果）に対しては小さな重み $g'(\tau) \ll 1$ しか乗らない} ことを意味する．
これはまさに「最悪の事態の発生確率を何倍にも大きく見積もって警戒する」という悲観的・リスク回避的な評価尺度を表している．
Moghimi and Ku \cite{moghimi2025beyond} のスペクトルリスク尺度（Spectral Risk Measure）$\int_0^1 \Finv(\tau)\phi(\tau)d\tau$ は，リスクスペクトル $\phi=g'$ と対応づければ式\eqref{eq:drm}と同等の枠組みである．

%--------------------------------------------------
\subsection{CVaR とその分位点表現}
\label{subsec:cvar}

\subsubsection{VaRの欠点とCVaRの導入}
金融実務における代表的なリスク指標として \textbf{VaR（Value at Risk; バリューアットリスク）} が広く知られている．
信頼水準 $1-\alpha$（例えば $95\%$，損失側の累積確率 $\alpha=0.05$）の VaR とは，$\alpha$-分位点 $\Finv(\alpha)$（の符号反転）である．
しかしVaR には実務上2つの課題が存在する．
\begin{enumerate}
  \item \textbf{裾の損失に対する情報欠落}：VaR は単なる境界値であり，その境界値より悪化方向の確率 $\alpha$ の領域における損失の大きさを評価できない．損失の大きさが異なっていても，境界値が同一であれば VaR の値は変化しない．
  \item \textbf{劣加法性の破綻}：分散投資によるリスク低減を反映する劣加法性 $\rho(X+Y) \le \rho(X)+\rho(Y)$ を，VaR は一般に満たさない（コヒーレントリスク尺度の公理である劣加法性を満たさない）．
\end{enumerate}

この欠点を克服するために考案されたのが \textbf{CVaR（Conditional Value at Risk; 条件付きバリューアットリスク，別名 Expected Shortfall）} である．

\begin{definition}[CVaR]
\label{def:cvar}
水準 $\alpha\in(0,1)$ に対し，損益 $Z$ の条件付きバリューアットリスクを
\begin{equation}
  \CVaR_\alpha(Z) = \frac{1}{\alpha}\int_0^\alpha\Finv(\tau)\,d\tau
  \label{eq:cvar}
\end{equation}
と定義する．
\end{definition}

式\eqref{eq:cvar}は，累積確率 $[0, \alpha]$（悪化方向の確率 $\alpha$ の領域）における分位点関数の平均値を表している．
% (簡潔化統合)
% (簡潔化統合)
収益分布が点質量をもたない連続分布である場合，これは条件付き期待値
\begin{equation}
  \CVaR_\alpha(Z) = \E\left[Z \;\middle|\; Z \le \Finv(\alpha)\right]
\end{equation}
と一致する．

\subsubsection{CVaRを導く歪み関数 $g_\alpha$}
式\eqref{eq:cvar}の CVaR は，定義\ref{def:drm}の歪みリスク尺度の特別な場合として自然に表せる．
歪み関数として
\begin{equation}
  g_\alpha(\tau) = \min\left(\frac{\tau}{\alpha},\,1\right)
  \label{eq:g-cvar}
\end{equation}
を考える．この関数の導関数を計算すると，
\begin{equation}
  g'_\alpha(\tau) = \begin{cases} \dfrac{1}{\alpha} & 0 \le \tau < \alpha \\[0.2cm] 0 & \alpha < \tau \le 1 \end{cases}
  \label{eq:gprime-cvar}
\end{equation}
となる．これを式\eqref{eq:drm}に代入すると，
\begin{equation}
  \rho_{g_\alpha}(Z) = \int_0^1 \Finv(\tau)\,g'_\alpha(\tau)\,d\tau = \int_0^\alpha \Finv(\tau)\,\frac{1}{\alpha}\,d\tau + \int_\alpha^1 \Finv(\tau)\cdot 0\,d\tau = \frac{1}{\alpha}\int_0^\alpha \Finv(\tau)\,d\tau = \CVaR_\alpha(Z)
\end{equation}
となり，式\eqref{eq:cvar}の定義と厳密に一致する．

この導関数 $g'_\alpha$ は，下側裾に重点を置く重み付けを表している．
具体的には，$[0, \alpha)$ の区間に $1/\alpha$ の重みを与え，$\alpha$ を超える領域の重みをゼロとする．
% (統合済み)
例えば $\alpha=0.05$ の場合，$[0, 0.05]$ の領域に一様の $1/0.05 = 20$ 倍の重みが与えられ，残りの $95\%$ の重みはゼロとなる．
この極端な重み付けこそが，第\ref{sec:intro}節で述べた「quantile fraction を下側裾 $[0, \alpha]$ に集中させるべきである」という直感の数学的裏付けである．

\subsubsection{Rockafellar--Uryasev 変分表現の完全解説}
Rockafellar and Uryasev \cite{rockafellar2000cvar} は，CVaR を以下の変分問題として定式化できることを示した．
\begin{equation}
  \CVaR_\alpha(Z) = \max_{\eta\in\R}\left\{\eta - \frac{1}{\alpha}\E\left[(\eta - Z)^+\right]\right\}.
  \label{eq:ru-cvar}
\end{equation}
ここで $(x)^+ = \max(x, 0)$ は正の部分（ReLU関数と同型）を表し，$\eta \in \R$ はしきい値を表す補助変数である（原著は損失の最小化として定式化しているが，本レポートでは利益に対する最大化として符号を反転している）．

この変分表現が $\CVaR_\alpha(Z)$ に一致することは，以下のように確認できる．
目的関数を $\psi(\eta) = \eta - \frac{1}{\alpha}\E[(\eta - Z)^+]$ とおく．
期待値を確率密度 $f(z)$ による積分で表すと，
\begin{equation}
  \psi(\eta) = \eta - \frac{1}{\alpha}\int_{-\infty}^\eta (\eta - z) f(z)\,dz
\end{equation}
である．これを $\eta$ で微分する．Leibnizの積分則を適用すると，
\begin{equation}
  \psi'(\eta) = 1 - \frac{1}{\alpha}\left( \underbrace{(\eta - \eta)f(\eta)}_{=\,0} + \int_{-\infty}^\eta 1\cdot f(z)\,dz \right) = 1 - \frac{F(\eta)}{\alpha}.
\end{equation}
停留条件 $\psi'(\eta) = 0$ を解くと，
\begin{equation}
  1 - \frac{F(\eta)}{\alpha} = 0 \iff F(\eta) = \alpha \iff \eta^* = \Finv(\alpha)
\end{equation}
となり，\textbf{最適解 $\eta^*$ は $\alpha$-分位点（VaR境界値）に一意に一致する}．
さらに2階微分は $\psi''(\eta) = -f(\eta)/\alpha \le 0$ であるから，$\psi$ は大域的凹関数であり，$\eta^*$ で最大値をとる．
この最適解 $\eta^* = \Finv(\alpha)$ を元の式に代入すると，
\begin{align}
  \psi(\eta^*)
  &= \Finv(\alpha) - \frac{1}{\alpha}\int_{-\infty}^{\Finv(\alpha)} (\Finv(\alpha) - z) f(z)\,dz \notag\\
  &= \Finv(\alpha) - \frac{1}{\alpha}\left( \Finv(\alpha)\int_{-\infty}^{\Finv(\alpha)} f(z)dz - \int_{-\infty}^{\Finv(\alpha)} z f(z)dz \right) \notag\\
  &= \Finv(\alpha) - \frac{1}{\alpha}\left( \Finv(\alpha)\cdot\alpha - \int_0^\alpha \Finv(\tau)\,d\tau \right) \notag\\
  &= \Finv(\alpha) - \Finv(\alpha) + \frac{1}{\alpha}\int_0^\alpha \Finv(\tau)\,d\tau = \CVaR_\alpha(Z)
\end{align}
となり，式\eqref{eq:cvar}の定義と一致することが示される．
この変分表現により，CVaR 最適化は凸計画問題に帰着されるため，ポートフォリオ最適化や機械学習において扱いやすいという利点をもつ．

%--------------------------------------------------
\subsection{歪み分布と推定誤差の上界}
\label{subsec:bound}

\subsubsection{歪み分布 $\mathcal{D}_g Z$ とは何か}
次に，本研究の理論的基礎となる \textbf{歪み分布} の概念を導入する．
$\hat Z$ をニューラルネットワーク等によって得られた $Z$ の近似分布（例えば式\eqref{eq:zhat_detailed}の離散分布 $\hat Z_N$）とする．
我々の究極の目的は，近似分布から計算したリスク尺度 $\rho_g(\hat Z)$ が，真のリスク尺度 $\rho_g(Z)$ とどれだけ近いかを厳密に評価することである．

\begin{definition}[歪み分布]
\label{def:distorted}
絶対連続な歪み関数 $g$ に対し，累積分布関数が $g$ である $[0,1]$ 上の確率変数 $U_g$ を考える（すなわち $\Prob(U_g \le u) = g(u)$）．
真の損益 $Z$ の分位点関数 $\Finv$ に対し，確率変数 $\Finv(U_g)$ が従う分布を \textbf{$Z$ の $g$ による歪み分布} と呼び，$\dist Z$（または $\mathcal{D}_g[Z]$）と表記する．
\end{definition}

この歪み分布 $\dist Z$ の性質を整理する．
一様乱数 $U \sim U(0, 1)$ を $\Finv$ に入力すると，逆関数法により元の確率変数 $Z = \Finv(U)$ が生成される．
歪み分布 $\dist Z$ とは，\textbf{「入力する乱数を一様分布から歪み分布 $U_g$ へと差し替えて生成される確率変数」} にほかならない．
その累積分布関数 $F_{\dist Z}(z)$ を計算すると，
\begin{equation}
  F_{\dist Z}(z) = \Prob\left(\Finv(U_g) \le z\right) = \Prob\left(U_g \le F_Z(z)\right) = g(F_Z(z)) = (g \circ F_Z)(z)
  \label{eq:dist-cdf}
\end{equation}
となる．元の累積分布関数 $F_Z(z)$ に歪み関数 $g$ を外側から合成した形となる．
さらに，確率密度関数 $f_{\dist Z}(z)$ を合成関数の微分で求めると，
\begin{equation}
  f_{\dist Z}(z) = \frac{d}{dz} g(F_Z(z)) = g'(F_Z(z))\,f_Z(z)
  \label{eq:dist-pdf}
\end{equation}
が得られる．
この式\eqref{eq:dist-pdf}は，歪み分布の確率密度関数を表している．
元の確率密度 $f_Z(z)$ に対し，その点が占める分位点位置の重み $g'(F_Z(z))$ が掛け算されて，密度が再重み付けされていることがわかる．
この歪み分布の期待値は次のように計算される．
\begin{equation}
  \E[\dist Z] = \int_{-\infty}^\infty z\,f_{\dist Z}(z)\,dz = \int_{-\infty}^\infty z\,g'(F_Z(z))f_Z(z)\,dz.
\end{equation}
ここで積分変数を分位レベル $\tau = F_Z(z)$（$d\tau = f_Z(z)dz$，$z = \Finv(\tau)$）に置換すると，
\begin{equation}
  \E[\dist Z] = \int_0^1 \Finv(\tau)\,g'(\tau)\,d\tau = \rho_g(Z)
  \label{eq:drm-as-mean}
\end{equation}
となる．
この結果から，歪みリスク尺度に関して次の関係が得られる．すなわち，歪みリスク尺度 $\rho_g(Z)$ は，歪み分布 $\dist Z$ における期待値 $\E[\dist Z]$ に他ならない．
% (地の文統合)
% (地の文統合)
% (地の文統合)
この関係により，歪みリスク尺度に対する解析を歪み分布に対する期待値の枠組みとして扱うことが可能となる．

\subsubsection{命題3.1（歪み重み付き $W_1$ 上界）の証明}
近似分布 $\hat Z$ を用いて計算した歪みリスク尺度の誤差 $|\rho_g(Z) - \rho_g(\hat Z)|$ の上界評価について述べる．
Pichler \cite{pichler2013evaluations} の確率測度の差異に対する評価に基づき，本研究の基本定理を以下に示す．

\begin{proposition}[歪み重み付き $\Wone$ 上界]
\label{prop:bound}
$g$ を絶対連続な歪み関数とし，$\Finv,\Fhinv$ は測度 $g'(\tau)d\tau$ に関して可積分であるとする．
このとき，任意の分布 $Z$ とその近似分布 $\hat Z$ に対し，
\begin{equation}
  \left|\rho_g(Z) - \rho_g(\hat Z)\right| \le \int_0^1 \left|\Finv(\tau) - \Fhinv(\tau)\right| g'(\tau)\,d\tau = \Wone(\dist Z, \dist\hat Z).
  \label{eq:bound}
\end{equation}
\end{proposition}

\begin{proof}
証明は2つのステップに分かれる．
第1のステップとして，左辺と中央の不等式を示す．
歪みリスク尺度の定義式\eqref{eq:drm}を直接代入すると，
\begin{align}
  \left|\rho_g(Z) - \rho_g(\hat Z)\right|
  &= \left| \int_0^1 \Finv(\tau)g'(\tau)d\tau - \int_0^1 \Fhinv(\tau)g'(\tau)d\tau \right| \notag\\
  &= \left| \int_0^1 \left(\Finv(\tau) - \Fhinv(\tau)\right) g'(\tau)\,d\tau \right|.
\end{align}
ここで積分の絶対値不等式（三角不等式 $\left|\int f\right| \le \int |f|$）を適用する．
歪み関数の単調性より導関数は常に非負（$g'(\tau) \ge 0$）であるから，
\begin{equation}
  \left| \int_0^1 \left(\Finv(\tau) - \Fhinv(\tau)\right) g'(\tau)\,d\tau \right| \le \int_0^1 \left|\Finv(\tau) - \Fhinv(\tau)\right| g'(\tau)\,d\tau
\end{equation}
となり，不等式が成り立つ．

第2のステップとして，中央の積分が右辺の Wasserstein 距離 $\Wone(\dist Z, \dist\hat Z)$ に一致することを示す．
1次元 1-Wasserstein 距離の分位点表現（式\eqref{eq:w1}）を歪み分布 $\dist Z$ と $\dist\hat Z$ に適用すると，
\begin{equation}
  \Wone(\dist Z, \dist\hat Z) = \int_0^1 \left| F_{\dist Z}^{-1}(v) - F_{\dist\hat Z}^{-1}(v) \right| dv.
\end{equation}
式\eqref{eq:dist-cdf}で見たように，歪み分布の累積分布関数は $F_{\dist Z} = g \circ F_Z$ である．
一般化逆関数をとると，$F_{\dist Z}^{-1}(v) = \Finv(g^{-1}(v))$ である．
したがって，
\begin{equation}
  \Wone(\dist Z, \dist\hat Z) = \int_0^1 \left| \Finv(g^{-1}(v)) - \Fhinv(g^{-1}(v)) \right| dv.
\end{equation}
ここで積分変数を $v$ から $\tau$ へと置換すると，
$v = g(\tau)$ とおくと，導関数から $dv = g'(\tau)d\tau$ であり，逆関数から $g^{-1}(v) = \tau$ である．
積分の範囲は，$v=0 \implies \tau = g^{-1}(0) = 0$，$v=1 \implies \tau = g^{-1}(1) = 1$ となる．
これらを代入すると，
\begin{equation}
  \int_0^1 \left| \Finv(g^{-1}(v)) - \Fhinv(g^{-1}(v)) \right| dv = \int_0^1 \left| \Finv(\tau) - \Fhinv(\tau) \right| g'(\tau)\,d\tau.
\end{equation}
これで第2の不等式も示された．
\end{proof}

\subsubsection{注意3.1：一様重みとの対比とLipschitz定数 $1/\alpha$}
\begin{remark}[なぜ一様重みではCVaR誤差の劣化が起きるのか]
\label{rem:bound-cvar}
命題\ref{prop:bound}の不等式は，本研究の目的関数設計の根拠を与える．
とくに CVaR 歪み $g_\alpha$ の場合，$g'_\alpha(\tau) = \frac{1}{\alpha}\mathbb{I}_{[0, \alpha)}(\tau)$ であるから，式\eqref{eq:bound}は
\begin{equation}
  \left|\CVaR_\alpha(Z) - \CVaR_\alpha(\hat Z)\right| \le \frac{1}{\alpha}\int_0^\alpha \left|\Finv(\tau) - \Fhinv(\tau)\right| d\tau
  \label{eq:bound-cvar}
\end{equation}
となる．
この式\eqref{eq:bound-cvar}から明らかなように，
右辺の積分区間は $[0, \alpha]$ のみであり，\textbf{$\alpha$ を超える領域 $[\alpha, 1]$（例えば $\alpha=0.05$ なら全体の $95\%$）における近似誤差は，CVaR の誤差上界に寄与しない}点である．

しかし，従来の FQF が最小化しているのは，一様重み 1-Wasserstein 距離
\begin{equation}
  \Wone(Z, \hat Z) = \int_0^\alpha \left|\Finv(\tau) - \Fhinv(\tau)\right| d\tau + \int_\alpha^1 \left|\Finv(\tau) - \Fhinv(\tau)\right| d\tau
\end{equation}
である．
FQF は，測度の上で大部分を占める $[\alpha, 1]$（長さ $1-\alpha = 0.95$ の広大な領域）の誤差を小さくしようと必死に quantile fraction を配分する．
すなわち，意思決定に関与しない領域（全体の $95\%$）の近似に表現リソースの大半が費やされていることになる．

さらに致命的なのは係数 $1/\alpha$ である．
Kiesel ら \cite{kiesel2016wasserstein} が指摘するように，CVaR（Expected Shortfall）は一様 1-Wasserstein 距離に関して\textbf{最悪 $1/\alpha$ の Lipschitz 定数}をもつ．すなわち，
\begin{equation}
  \left|\CVaR_\alpha(Z) - \CVaR_\alpha(\hat Z)\right| \le \frac{1}{\alpha}\,\Wone(Z, \hat Z).
\end{equation}
もし同じ $\Wone(Z, \hat Z) = 0.1$ という小さな誤差を達成していたとしても，その誤差がたまたま下側裾 $[0, \alpha]$ に集中していた場合，CVaR の推定誤差は最大で $0.1 \times (1/0.05) = 2.0$ へと \textbf{最大 20倍に拡大される}ことになる．
一様重みでの近似に甘んじている限り，この $1/\alpha$ 倍の増幅爆弾を回避することは原理的に不可能である．
これが，「FQF の一様重み $\Wone$ を歪み重み付き $\Wone(\dist Z, \dist\hat Z)$ に置き換える」という本研究の動機である．
\end{remark}

%--------------------------------------------------
\subsection{離散近似への適用と正則化を含まない $g_\alpha$ による評価}
\label{subsec:discrete-eval}

\subsubsection{離散分布に対する閉形式計算}
第\ref{subsec:fqf}節の式\eqref{eq:zhat_detailed}で定義した，FQF 型の離散近似分布
\begin{equation}
  \hat Z_N = \sum_{i=0}^{N-1}(\tau_{i+1}-\tau_i)\,\delta_{\Finv(\hat\tau_i)}
\end{equation}
これに対して，歪みリスク尺度 $\rho_g(\hat Z_N)$ は数値積分を要さず，次の閉形式で計算できる．
\begin{equation}
  \rho_g(\hat Z_N) = \sum_{i=0}^{N-1}\left(g(\tau_{i+1}) - g(\tau_i)\right)\Finv(\hat\tau_i).
  \label{eq:rho-discrete}
\end{equation}
近似分布 $\hat Z_N$ の分位点関数 $\Fhinv(u)$ は，各区間 $u \in (\tau_i, \tau_{i+1}]$ で定数値 $\Finv(\hat\tau_i)$ をとる階段関数であるから，
\begin{align}
  \rho_g(\hat Z_N)
  &= \int_0^1 \Fhinv(u)\,g'(u)\,du = \sum_{i=0}^{N-1} \int_{\tau_i}^{\tau_{i+1}} \Finv(\hat\tau_i)\,g'(u)\,du \notag\\
  &= \sum_{i=0}^{N-1} \Finv(\hat\tau_i) \int_{\tau_i}^{\tau_{i+1}} g'(u)\,du = \sum_{i=0}^{N-1} \left(g(\tau_{i+1}) - g(\tau_i)\right)\Finv(\hat\tau_i)
\end{align}
となり，式\eqref{eq:rho-discrete}が得られる．

\paragraph{分位点交差（quantile crossing）に関する注意}
式\eqref{eq:rho-discrete}が離散確率分布 $\hat Z_N$ の歪みリスク尺度（特に CVaR）と一致するのは，各原子の値が単調増加（分位点の交差がない）している場合に限る．
\begin{equation*}
  \Finv(\hat\tau_0) \le \Finv(\hat\tau_1) \le \dots \le \Finv(\hat\tau_{N-1})
\end{equation*}
を満たしている（分位点交差がない）場合に限る．真の分位点関数では単調性が成り立つため実験Aでは厳密に成立する．しかし実験Bのように深層ニューラルネットワークで分位点値を出力する場合，ネットワークの局所的な近似誤差により分位点の逆転（quantile crossing）が生じうる．交差が生じた場合，式\eqref{eq:rho-discrete}は厳密な離散分布の CVaR ではなく，「ネットワーク出力が（単調性の崩れた）近似分位点関数に対する数値積分（quadrature）値」として機能する．厳密な離散分布としての CVaR を評価するには，評価時に原子の値を昇順にソートし対応する累積重みを割り当て直す（あるいは経験分布を用いる）必要があるが，本レポートでは FQF や IQN の標準的慣例に従い，分位点関数出力に対する線形結合値としての位置づけを一貫して採用する（なお実験Bの学習済みモデルにおいて主要な意思決定領域で顕著な交差は観測されていない）．

\subsubsection{正則化歪み $g_{\alpha,\lambda}$ の最適化と純粋 CVaR 評価の理論的整合性}
第\ref{sec:proposed}節で後述するように，実際のアルゴリズム最適化においては，勾配消失を防ぐために微小な正則化項 $\lambda > 0$ を付加した正則化歪み関数
\begin{equation}
  g_{\alpha,\lambda}(\tau) = (1-\lambda)\,g_\alpha(\tau) + \lambda\tau
\end{equation}
（既定値 $\lambda=0.05$）を用いる．
ここで，最適化に用いる目的関数と最終的に評価する CVaR（$\lambda=0$）の差異について検討する．
評価時には，式\eqref{eq:rho-discrete}に正則化を含まない純粋な $g_\alpha$ を代入して $\CVaR_\alpha(\hat Z_N)$ を計算する．

このとき，正則化目的 $L_{g_{\alpha,\lambda}}$ を最小化することは，純粋な CVaR 重み付き誤差
\begin{equation}
  L_{g_\alpha} = \frac{1}{\alpha}\int_0^\alpha \left|\Finv(u) - \Fhinv(u)\right|du
\end{equation}
本研究において，両者の間には以下の不等式が成立する．
% 不等式の成立
導関数 $g'_{\alpha,\lambda}(u)$ は，$[0, \alpha)$ において
\begin{equation}
  g'_{\alpha,\lambda}(u) = (1-\lambda)\,\frac{1}{\alpha} + \lambda = \frac{1-\lambda+\lambda\alpha}{\alpha} = (1-\lambda+\lambda\alpha)\,g'_\alpha(u)
\end{equation}
と定数倍の関係にある．
したがって，$[0, \alpha]$ の範囲において，被積分関数は $g'_\alpha(u) = \frac{1}{1-\lambda+\lambda\alpha} g'_{\alpha,\lambda}(u)$ と厳密に書き換えられる．
これより，
\begin{align}
  L_{g_\alpha}
  &= \int_0^\alpha \left|\Finv(u) - \Fhinv(u)\right| g'_\alpha(u)\,du \notag\\
  &= \frac{1}{1-\lambda+\lambda\alpha}\int_0^\alpha \left|\Finv(u) - \Fhinv(u)\right| g'_{\alpha,\lambda}(u)\,du \notag\\
  &\le \frac{1}{1-\lambda+\lambda\alpha}\int_0^1 \left|\Finv(u) - \Fhinv(u)\right| g'_{\alpha,\lambda}(u)\,du = \frac{1}{1-\lambda+\lambda\alpha}\,L_{g_{\alpha,\lambda}}
  \label{eq:reg-bound-relation}
\end{align}
という厳密な不等式が得られる（最後の不等号は，被積分関数が非負であるため積分区間を $[0, \alpha]$ から $[0, 1]$ へ広げたことによる）．

基本設定値である $\lambda=0.05$ および $\alpha=0.05$ を代入すると，次式が得られる．
\begin{equation}
  \frac{1}{1-\lambda+\lambda\alpha} = \frac{1}{1 - 0.05 + 0.05\times0.05} = \frac{1}{0.9525} \approx 1.04987.
\end{equation}
すなわち，\textbf{正則化目的関数 $L_{g_{\alpha,\lambda}}$ を最小化することは，評価対象である純粋 CVaR 誤差 $L_{g_\alpha}$ を高々約 5\% の係数差で同時に抑え込むことを理論的に理論的に保証している}．
この数理的保証があるからこそ，我々は正則化歪み $g_{\alpha,\lambda}$ の導入が理論的に正当化される．

%==================== 4 提案モデル ====================
\section{提案モデル：歪み重み付きquantile fraction提案}
\label{sec:proposed}

\subsection{FPN の損失の置換}
前章までの準備に基づき，本章では提案手法である歪み FQF について述べる．
従来の FQF（第\ref{subsec:fqf}節）において，Fraction Proposal Network（FPN）が最小化していた母集団目的関数は，一様重み 1-Wasserstein 距離 $\Wone(Z, \hat Z_N)$ であった．
本研究では，この目的関数を，命題\ref{prop:bound}で導出した歪み重み付き 1-Wasserstein 距離
\begin{equation}
  L_g(\tau) = \int_0^1 \left|\Finv(u) - \Fhinv(u)\right| g'(u)\,du = \Wone(\dist Z, \dist\hat Z_N)
  \label{eq:Lg}
\end{equation}
に置き換える．ここで $\tau = (\tau_1, \dots, \tau_{N-1})$ は FPN が出力する内部の quantile fraction 列である．
命題\ref{prop:bound}により，この $L_g(\tau)$ を最小化することは，歪みリスク尺度 $\rho_g$ の推定誤差の上界を直接的に最小化することと同値である．
一様歪み $g(\tau)=\tau$（$g'(\tau)=1$）のときは $L_g = \Wone(Z, \hat Z_N)$ となり従来 FQF に一致するため，提案手法は従来 FQF を特別な場合として内包する一般化となっている．

%--------------------------------------------------
\subsection{最適quantile fraction＝歪み分布の1次元最適量子化}

\subsubsection{命題4.1（$g$-中央値原子）の導出}
目的関数を歪み重み付き $L_g$ に拡張したとき，各区間 $[\tau_i, \tau_{i+1}]$ を代表する原子の位置 $q_i$ の最適配置を導出する．
一様重みのときは算術中点 $\Finv((\tau_i+\tau_{i+1})/2)$ であったが，歪み重み $g'(u)$ の下では重みに応じて最適点がシフトする．

\begin{proposition}[$g$-中央値原子]
\label{prop:gmedian}
quantile fraction $\tau_0 < \dots < \tau_N$ を固定する．
区間 $i$ の目的関数への寄与
\begin{equation}
  \ell_i^g(q) = \int_{\tau_i}^{\tau_{i+1}} \left|\Finv(u) - q\right| g'(u)\,du
  \label{eq:g-loss-cell}
\end{equation}
を最小化する最適原子 $q_i$ は，
\begin{equation}
  q_i = \Finv(\hat\tau_i),\qquad \hat\tau_i = g^{-1}\!\left(\frac{g(\tau_i) + g(\tau_{i+1})}{2}\right)
  \label{eq:gmedian}
\end{equation}
で与えられる．すなわち，$\hat\tau_i$ は区間上の測度 $g'(u)du$ に関する中央値（$g$-中央値）である．
\end{proposition}

\begin{proof}
第\ref{subsec:fqf}節の補題\ref{lem:fqf-atom}と同様に，行間を埋めて証明する．
$\Finv$ は連続かつ非減少であるから，$q = \Finv(m)$ を満たす点 $m \in [\tau_i, \tau_{i+1}]$ が存在する．
積分区間を $m$ で分割し，絶対値を外すと，
\begin{equation}
  \ell_i^g(q) = \int_{\tau_i}^{m(q)} \left(q - \Finv(u)\right) g'(u)\,du + \int_{m(q)}^{\tau_{i+1}} \left(\Finv(u) - q\right) g'(u)\,du.
\end{equation}
これを $q$ で微分する．Leibnizの積分則により，端点 $m(q)$ の微分項は，被積分関数が $u=m$ において
$(q - \Finv(m))g'(m) = 0$ および $(\Finv(m) - q)g'(m) = 0$
となるため，厳密に消滅する．
残る被積分関数の $q$ 微分を計算すると，
\begin{align}
  \frac{\partial \ell_i^g}{\partial q}
  &= \int_{\tau_i}^m 1\cdot g'(u)\,du - \int_m^{\tau_{i+1}} 1\cdot g'(u)\,du \notag\\
  &= \left[g(u)\right]_{\tau_i}^m - \left[g(u)\right]_m^{\tau_{i+1}} \notag\\
  &= \left(g(m) - g(\tau_i)\right) - \left(g(\tau_{i+1}) - g(m)\right) \notag\\
  &= 2g(m) - \left(g(\tau_i) + g(\tau_{i+1})\right).
  \label{eq:gmedian-deriv}
\end{align}
最小点における停留条件 $\frac{\partial \ell_i^g}{\partial q} = 0$ より，
\begin{equation}
  2g(m) - \left(g(\tau_i) + g(\tau_{i+1})\right) = 0 \iff g(m) = \frac{g(\tau_i) + g(\tau_{i+1})}{2}.
\end{equation}
両辺に一般化逆関数 $g^{-1}$ を作用させると，
\begin{equation}
  m = g^{-1}\!\left(\frac{g(\tau_i) + g(\tau_{i+1})}{2}\right) = \hat\tau_i
\end{equation}
が得られる．2階微分は $\frac{\partial^2 \ell_i^g}{\partial q^2} > 0$ であるから，これが大域的最小点である．
\end{proof}

式\eqref{eq:gmedian}は次のような幾何学的意味をもつ．
式\eqref{eq:gmedian-deriv}の第一項 $g(m) - g(\tau_i)$ は点 $m$ の左側の歪み測度質量であり，第二項 $g(\tau_{i+1}) - g(m)$ は右側の歪み測度質量である．
それが等しく釣り合う点，すなわち \textbf{「区間上の測度 $g'(u)du$ の質量をちょうど半分ずつに二等分する中央値（$g$-中央値）」} における分位点値こそが，誤差を最小化する最適な原子の位置に他ならない．

\paragraph{各区間における $g$-中央値の振る舞い}
正則化歪み $g_{\alpha,\lambda}$ のもとで，この $g$-中央値 $\hat\tau_i$ は区間 $[\tau_i, \tau_{i+1}]$ の位置に応じて以下の3つの場合に分類される．
% (重複行の削除)
\begin{enumerate}
  \item \textbf{区間が下側裾にある場合（$\tau_{i+1} \le \alpha$）}：
  この領域では $g_{\alpha,\lambda}(\tau)$ は傾き $(1-\lambda)/\alpha + \lambda$ の純粋な線形関数（1次式）である．
  1次関数の下では中点の像は像の中点と一致するため，
  $g$-中央値は通常の算術中点 $\hat\tau_i = (\tau_i + \tau_{i+1})/2$ と厳密に一致する．
  \item \textbf{区間が裾野外（$\tau_i \ge \alpha$）にある場合}：
  この領域でも $g_{\alpha,\lambda}(\tau)$ は傾き $\lambda$ の線形関数であるため，同様に算術中点 $\hat\tau_i = (\tau_i + \tau_{i+1})/2$ と一致する．
  \item \textbf{区間が境界 $\alpha$ を跨ぐ場合（$\tau_i < \alpha < \tau_{i+1}$）}：
  この区間でのみ折れ線の効果が現れる．$[0, \alpha]$ 側の急な傾き（大きな重み）に引っ張られるため，
  $g$-中央値 $\hat\tau_i$ は通常の算術中点よりも左側（下側裾の損失側）へとシフトする．
\end{enumerate}
すなわち，裾野内部では算術中点により二次精度（注意\ref{rem:midpoint}）を保ちつつ，境界付近でのみ歪み密度に応じて適応的シフトが行われる実用的な性質をもつ．


\subsubsection{最適量子化（Optimal Quantization）としての数学的同値性}
命題\ref{prop:bound}（誤差上界）と命題\ref{prop:gmedian}（$g$-中央値原子）を組み合わせることで，目的関数 $L_g$ の最小化は，\textbf{最適量子化（Optimal Quantization）} \cite{graf2000quantization} に帰着される．

累積確率軸に対して，変数変換 $v = g(\tau)$（$dv = g'(\tau)d\tau$）を適用する．
すると，累積確率軸 $[0, 1]$ は歪み空間 $[0, 1]$ へと写像され，
各境界点 $\tau_i$ は $v_i = g(\tau_i)$ へ，各原子の代表レベル $\hat\tau_i$ は
\begin{equation}
  \hat v_i = g(\hat\tau_i) = \frac{g(\tau_i) + g(\tau_{i+1})}{2} = \frac{v_i + v_{i+1}}{2}
\end{equation}
となり，歪み空間上での算術中点へと写像される．
そして目的関数式\eqref{eq:Lg}は，
\begin{equation}
  L_g(\tau) = \int_0^1 \left| F_{\dist Z}^{-1}(v) - F_{\dist\hat Z_N}^{-1}(v) \right| dv = \Wone(\dist Z, \dist\hat Z_N)
\end{equation}
と表される．
これは，Graf and Luschgy \cite{graf2000quantization} が体系化した，
\textbf{「歪み確率分布 $\dist Z$ を，1次元空間において $N$ 個の離散原子によって $L^1$ ノルムの意味で最も忠実に近似する 1次元最適量子化問題（$r=1$）」}
そのものである．

1次元空間において，任意の離散点集合に対する Voronoi 領域（各点に最も近い空間の領域）は必ず「区間（interval）」になるため，累積確率軸を区間分割する本パラメータ化は一般性を失わない．
最適量子化の理論からは，以下の性質が得られる．
\begin{itemize}
  \item \textbf{Lloyd 型の固定点条件}：最適量子化器は，「原子＝各区間の $g$-中央値」「境界＝隣接原子の中点」という連立固定点条件を満たす．
  \item \textbf{Zador（ザドール）型定理}：適切な正則性条件下で，$N$ 点最適量子化誤差の漸近オーダーは $O(1/N)$ で減衰する．
\end{itemize}

\subsubsection{点密度比 $1/\alpha$ と改善の幾何学的直観}
CVaR の歪み $g_\alpha$ のもとでの最適量子化の挙動について考察する．
歪み測度 $g'_\alpha$ は $[0, \alpha]$ にのみ集中しており，$\alpha$ より上の領域には質量が存在しない（測度ゼロ）．
したがって，最適量子化器は，与えられた $N$ 個の表現リソース（原子）の \textbf{すべてを下側裾 $[0, \alpha]$ に集中して配置する}．
これに対し，従来の FQF や QR-DQN のような一様配置が同じ領域 $[0, \alpha]$ に割り当てる点数は，全体の $\alpha$ 割（わずか $\alpha N$ 個）にすぎない．

したがって，下側裾 $[0, \alpha]$ における点の密度比は，
\begin{equation}
  \frac{N}{\alpha N} = \frac{1}{\alpha}
\end{equation}
となる．$\alpha=0.05$ であれば，点密度は\textbf{20倍} に達する．
この密度の圧倒的な差が，CVaR 推定精度の向上をもたらす定量的要因である．

\subsubsection{予想4.1（改善スケーリング）とヘビーテイルにおける端点特異性}
この直感を，裾野の発散度合い（正則変動指数）まで考慮して厳密に定式化したのが次の予想である．

\begin{conjecture}[誤差上界の比による改善スケーリング]
\label{conj:alpha}
$\Finv$ が $(0,\alpha]$ で $C^1$ 級かつ $\Finv\in L^1[0,\alpha]$（$\CVaR_\alpha$ が有限値をもつ）であり，さらに導関数が $(\Finv)'(\tau)=\tau^{-\beta}\ell(\tau)$ の形で原点近傍（$\tau\to+0$）において正則変動する（regular variation）分布族を考える．
ここで $\beta$ は原点付近における分位点関数の発散の強さを表す指数（自由度 $\nu$ の $t$ 分布では $\beta=1+1/\nu$，ガウス裾では $\beta=1$），$\ell(\tau)$ は対数のように多項式より緩やかに変化する関数（緩変動関数）である．
このとき，$\alpha N\ge2$ の領域において，一様quantile fractionによる式\eqref{eq:bound-cvar}の上界の下界 $E_{\mathrm{unif}}$ と，歪み最適quantile fractionによる上界 $E_{\mathrm{opt}}$ の比は，$\beta\le1$（ガウス裾など）のとき漸近オーダーで $\Theta(1/\alpha)$，$\beta\in(1,2)$（自由度 $\nu$ の $t$ 分布などのヘビーテイル分布）のとき
\begin{equation}
  \frac{E_{\mathrm{unif}}}{E_{\mathrm{opt}}} \;\asymp\; \frac{N^{\beta-1}}{\alpha^{2-\beta}}
  \label{eq:conj-scaling}
\end{equation}
型となる（ここで記号 $\asymp$ は，両辺の比が正の上下限で挟まれ，漸近的に同じオーダーで変化することを表す）．
とくに $\alpha N\ge C$（$C$ は定数）の領域では，いずれの場合もこの改善比は少なくとも $\Omega(1/\alpha)$ となる．
\end{conjecture}

\paragraph{予想\ref{conj:alpha}における2つの要因}
式\eqref{eq:conj-scaling}が示す数理的メカニズムは，以下の2つの効果に分解して理解できる．
改善比は，以下の2つの効果の積として整理できる．
\begin{enumerate}
  \item \textbf{集中配置による $1/\alpha$ 効果（分母）}：
  裾野 $[0, \alpha]$ への点の集中配置により，点密度が $1/\alpha$ 倍になる効果である．
  \item \textbf{ヘビーテイルにおける端点特異性の自律的回避効果（分子 $N^{\beta-1}$）}：
  $t$分布のように裾の重いヘビーテイル分布では，$\tau \to 0$ において分位点関数がマイナス無限大に向かって急速に発散する（特異性）．
  一様配置では，第1区間 $[0, 1/N]$ がこの特異性を丸ごと抱え込んでしまうため，誤差の減衰速度が $N^{-1}$ よりも大幅に遅い $O(N^{-(2-\beta)}) = O(N^{-(1-1/\nu)})$（$\nu=3$ なら $N^{-2/3}$）へと致命的に劣化する．
  これに対し，歪み最適配置は，後述の変分原理（注意\ref{rem:midpoint}）により，特異点に向かって区間幅を自律的に極小化（$\tau_1 \propto N^{-3}$）し，特異性による減衰劣化を回避できる．
\end{enumerate}
この「減衰指数の差（$N^{\beta-1}$, $\beta > 1$）」が分子に現れるため，ヘビーテイルな環境であればあるほど，そして原子数 $N$ を増やすほど，提案手法の一様配置に対する改善比は $1/\alpha$ を超えて拡大する．

\subsubsection{注意4.1：変分法による点密度の導出と中点則二次精度}
\begin{remark}[変分法による最適点密度定理と中点則]
\label{rem:midpoint}
正則化歪み $g_{\alpha,\lambda}$ は $[0,\alpha)$ 上で線形であるため，$g$-中央値は算術中点と一致し，歪み最適配置による CVaR 推定は $[0,\alpha]$ 上の\textbf{非一様格子}中点則に一致する．
歪み最適配置によって特異性の影響が抑えられる理由について，変分法（Euler--Lagrange 方程式）を用いて述べる．

正則な領域において，区間幅 $w_i = \tau_{i+1} - \tau_i$ をもつ区間 $i$ の中点則による近似誤差は，Taylor展開の2次項から
\begin{equation}
  \int_{\tau_i}^{\tau_{i+1}} |\Finv(u) - \Finv(\hat\tau_i)|\,du \approx \frac{1}{4}\,(\Finv)'(\tau_i)\,w_i^2
\end{equation}
と評価できる．
ここで，$[0, \alpha]$ 上に配置される quantile fraction の連続的な「点密度関数」を $p(\tau)$ と置く（局所的な区間幅は $w_i \approx \frac{1}{N p(\tau_i)}$ と反比例する）．
全区間にわたる総誤差は，リーマン和から積分へと移行して，
\begin{equation}
  E_{\mathrm{total}} \approx \sum_{i} \frac{1}{4}(\Finv)'(\tau_i)\left(\frac{1}{N p(\tau_i)}\right)^2 \approx \frac{1}{4N}\int_0^\alpha \frac{(\Finv)'(\tau)}{p(\tau)}\,d\tau
  \label{eq:integral-error}
\end{equation}
と近似できる．
我々の目標は，総点数の制約条件 $\int_0^\alpha p(\tau)\,d\tau = 1$（$p(\tau) \ge 0$）のもとで，総誤差式\eqref{eq:integral-error}を最小化する最適な点密度 $p(\tau)$ を見つけ出すことである．

これは汎関数の変分問題である．
ラグランジュ未定乗数 $\mu > 0$ を導入し，作用積分を
\begin{equation}
  \mathcal{J}[p] = \int_0^\alpha \left[ \frac{(\Finv)'(\tau)}{p(\tau)} + \mu\,p(\tau) \right] d\tau - \mu
\end{equation}
と定義する．被積分関数を $\mathcal{F}(p) = \frac{(\Finv)'(\tau)}{p} + \mu p$ とおき，$p$ で変分をとると，
\begin{equation}
  \frac{\partial\mathcal{F}}{\partial p} = -\frac{(\Finv)'(\tau)}{p(\tau)^2} + \mu = 0 \iff p(\tau)^2 = \frac{1}{\mu}\,(\Finv)'(\tau).
\end{equation}
両辺の正の平方根をとることで，
\begin{equation}
  p(\tau) \propto \sqrt{(\Finv)'(\tau)}
  \label{eq:density-theorem}
\end{equation}
となり，\textbf{「1次元 $L^1$ 最適量子化の点密度定理（Graf and Luschgy 2000）」} に帰着される．
この変分解により，下側裾においてより密に原子が配置される．

自由度 $\nu$ の $t$ 分布では $(\Finv)'(\tau) \sim \tau^{-(1+1/\nu)}$ であるから，点密度は $p(\tau) \propto \tau^{-(1+1/\nu)/2}$ となる．
第1区間の条件 $\int_0^{\tau_1} p(\tau)d\tau = 1/N$ を積分計算すると，第1区間幅は $\tau_1 \propto N^{-2\nu/(\nu-1)}$（$\nu=3$ では $\tau_1 \propto N^{-3}$）と極限的なスピードで縮小する．
このとき，第1区間の端点特異性の寄与は $\tau_1^{1-1/\nu} \propto (N^{-3})^{2/3} = N^{-2}$ となり，自由度 $\nu$（ヘビーテイルの重さ）に依存する減衰劣化が解消される．
したがって，歪み最適配置では二次精度 $O(N^{-2})$ の収束が得られる．
\end{remark}

%--------------------------------------------------
\subsection{quantile fractionに関する境界勾配の導出}
\label{subsec:dist-grad}

\subsubsection{命題4.2（提案歪み境界勾配）の導出}
FPN のパラメータを勾配降下法で学習させるためには，目的関数 $L_g$ の各境界 $\tau_i$（$i=1,\dots,N-1$）に関する勾配が必要である．
この命題は，提案歪み境界勾配が従来 FQF の境界勾配式\eqref{eq:fqf-prop1}に重み係数 $g'(\tau_i)$ を乗じた簡潔な閉形式になることを示している．

\begin{proposition}[歪み重み付き $\Wone$ に対する境界勾配]
\label{prop:grad}
$\Finv$ は連続かつ非減少であり，境界 $\tau_i$ および原子 $q_i = \Finv(\hat\tau_i)$ が式\eqref{eq:gmedian}の $g$-中央値に置かれているとする．さらに各 $\tau_i$ は $g$ の微分可能点にある（$\tau_i \ne \alpha$）とする．このとき $i=1,\dots,N-1$ について
\begin{equation}
  \frac{\partial L_g}{\partial\tau_i} = g'(\tau_i)\left[2\Finv(\tau_i) - \Finv(\hat\tau_i) - \Finv(\hat\tau_{i-1})\right].
  \label{eq:dist-grad}
\end{equation}
\end{proposition}

\begin{proof}
目的関数の総和式 $L_g = \sum_{j=0}^{N-1} \ell_j^g(q_j)$ において，境界 $\tau_i$ に依存するのは区間 $i-1$（上限が $\tau_i$）と区間 $i$（下限が $\tau_i$）の2つのみである．
第\ref{subsec:fqf}節の命題\ref{prop:fqf-grad}と同様に，勾配の展開を行う．
\begin{equation}
  \frac{\partial L_g}{\partial\tau_i} = \frac{\partial \ell_{i-1}^g}{\partial\tau_i}\Bigg|_{q_{i-1}\text{固定}} + \frac{\partial \ell_i^g}{\partial\tau_i}\Bigg|_{q_i\text{固定}} + \underbrace{\frac{\partial \ell_{i-1}^g}{\partial q_{i-1}}\,\frac{\partial q_{i-1}}{\partial\tau_i} + \frac{\partial \ell_i^g}{\partial q_i}\,\frac{\partial q_i}{\partial\tau_i}}_{\text{包絡線定理によりゼロ}}.
\end{equation}
各原子 $q_j$ は命題\ref{prop:gmedian}により区間損失 $\ell_j^g$ を最小化するように選ばれているため，
極値条件 $\frac{\partial \ell_{i-1}^g}{\partial q_{i-1}} = 0$ および $\frac{\partial \ell_i^g}{\partial q_i} = 0$ が恒等的に成り立つ．
したがって，原子の位置を通じた間接的な依存性は厳密にゼロとなる（包絡線定理）．

残る端点の直接微分を Leibniz の積分則により計算する．
\begin{itemize}
  \item \textbf{区間 $i-1$ の積分上限の微分}（$\int_{\tau_{i-1}}^{\tau_i} |\Finv(u) - q_{i-1}| g'(u)du$）：
  $+|\Finv(\tau_i) - q_{i-1}|\,g'(\tau_i)$．
  \item \textbf{区間 $i$ の積分下限の微分}（$\int_{\tau_i}^{\tau_{i+1}} |\Finv(u) - q_i| g'(u)du$）：
  $-|\Finv(\tau_i) - q_i|\,g'(\tau_i)$．
\end{itemize}
したがって，
\begin{equation}
  \frac{\partial L_g}{\partial\tau_i} = g'(\tau_i)\left( |\Finv(\tau_i) - q_{i-1}| - |\Finv(\tau_i) - q_i| \right).
\end{equation}
単調性 $\Finv(\hat\tau_{i-1}) \le \Finv(\tau_i) \le \Finv(\hat\tau_i)$ により絶対値を外すと，
\begin{align}
  \frac{\partial L_g}{\partial\tau_i}
  &= g'(\tau_i)\left[ \left(\Finv(\tau_i) - \Finv(\hat\tau_{i-1})\right) - \left(\Finv(\hat\tau_i) - \Finv(\tau_i)\right) \right] \notag\\
  &= g'(\tau_i)\left[ 2\Finv(\tau_i) - \Finv(\hat\tau_i) - \Finv(\hat\tau_{i-1}) \right].
\end{align}
以上により証明される．
\end{proof}

\subsubsection{実装上の勾配と代理損失（stop-gradient 構造）}
ここで，深層学習（PyTorch等）で本手法を実装する際に，深層強化学習の実装において留意すべき\textbf{勾配計算と代理損失の設計}について述べる．

数式\eqref{eq:dist-grad}の勾配値
\begin{equation}
  d_i = g'(\tau_i)\left[2\Finv_{\omega_2}(\tau_i) - \Finv_{\omega_2}(\hat\tau_i) - \Finv_{\omega_2}(\hat\tau_{i-1})\right]
  \label{eq:grad-val}
\end{equation}
この FPN の逆伝播において，目的関数 $L_g$ に対して PyTorch の自動微分（\texttt{loss.backward()}）を直接適用することには留意が必要である．
通常の自動微分をそのまま適用した場合，以下の2つの問題が生じる．
\begin{enumerate}
  \item \textbf{分位点値ネットワーク $\omega_2$ への勾配の流出}：
  分位点値ネットワーク $\Finv_{\omega_2}$ は Huber 分位点回帰損失（式\eqref{eq:qr-loss}）によってのみ学習されるべきである．もし $L_g$ の微分が $\omega_2$ の重みにまで逆伝播してしまうと，分位点値ネットワークが「FPN の都合の良いように値を歪める」という干渉を起こし，学習が崩壊する．
  \item \textbf{意図しない高階微分の混入}：
  式\eqref{eq:grad-val}の係数 $g'(\tau_i)$ 自体も $\tau_i$ の関数である．自動微分を行うと，連鎖律によって $\tau_i$ が $g'$ を微分してしまい，理論には存在しない $g''(\tau_i)$ の項（高階微分）が勝手に勾配に紛れ込んでしまう．
\end{enumerate}

意図するのは，$\Finv_{\omega_2}$ や $g'$ への微分を遮断した上で，数値ベクトルとしての勾配値 $d_i$ を算出し，これを FPN の出力 $\tau_i$ のみに逆伝播させることである．
実装（\texttt{fqf.py}）では，代理損失（Surrogate Loss）と stop-gradient（勾配停止）の技法を採用している．

\begin{enumerate}
  \item 勾配値 $d_i$（$i=1,\dots,N-1$）の計算は，勾配計算を遮断したコンテキスト
  \begin{center}
    \texttt{with torch.no\_grad():}
  \end{center}
  の内部でテンソルとして計算する（$\operatorname{stopgrad}(d_i)$）．
  \item FPN が出力した各予測境界 $\tau_i$ に対して，次のような線形な代理損失を定義する．
  \begin{equation}
    L_{\mathrm{frac}} = \sum_{i=1}^{N-1} \operatorname{stopgrad}(d_i)\,\tau_i.
    \label{eq:surrogate-loss}
  \end{equation}
  \item この代理損失 $L_{\mathrm{frac}}$ に対して \texttt{L\_frac.backward()} を実行する．
\end{enumerate}
代理損失\eqref{eq:surrogate-loss}を $\tau_i$ で偏微分すると，定義から直ちに
\begin{equation}
  \frac{\partial L_{\mathrm{frac}}}{\partial \tau_i} = d_i
\end{equation}
となる．
こうして，何ら余計な副次効果を生むことなく，理論通りの勾配ベクトル $d_i$ が FPN のパラメータ $\omega_1$ にのみ正確に逆伝播されるのである．
これは FQF の公式実装でも採用されている深層強化学習における標準的な設計である．
深層学習フレームワーク（PyTorch）における具体的なコード構成を以下に示す．
\begin{quote}
\small
\begin{verbatim}
# 1. 勾配値 d_i の計算（勾配追跡を遮断）
with torch.no_grad():
    # 命題4.2の理論勾配ベクトル: g'(tau) * (2*Q(tau) - Q(hat_tau) - Q(prev_hat_tau))
    d = g_prime(tau[:, 1:-1]) * (2 * Q_tau - Q_hat_tau - Q_prev_hat_tau)

# 2. 代理損失の定義（FPNの出力 tau に対する線形和）
surrogate_loss = (d * tau[:, 1:-1]).sum(dim=-1).mean()

# 3. FPN のみに勾配を逆伝播
surrogate_loss.backward()
\end{verbatim}
\end{quote}
このように，包絡線定理に基づいた FPN への勾配計算は簡潔に実装できる．


%--------------------------------------------------
\subsection{正則化歪み関数}
\label{subsec:g-reg}

\subsubsection{純粋な $g_\alpha$ を用いた場合の課題}
ここで，CVaR の歪み関数 $g_\alpha$ を直接最適化に用いず，正則化歪み関数 $g_{\alpha,\lambda}$ を導入する理由について述べる．
% 純粋なCVaR歪みの課題

純粋な CVaR 歪み $g_\alpha$ の導関数は，式\eqref{eq:gprime-cvar}で見たように
\begin{equation*}
  g'_\alpha(\tau) = \begin{cases} 1/\alpha & 0 \le \tau < \alpha \\ 0 & \alpha < \tau \le 1 \end{cases}
\end{equation*}
これをそのまま式\eqref{eq:dist-grad}の勾配計算に用いると，以下の2つの問題が生じる．
\begin{enumerate}
  \item \textbf{勾配消失問題（死んだ境界点）}：
  $\tau_i > \alpha$ に位置するすべての境界点において，$g'(\tau_i) = 0$ となるため勾配が厳密にゼロとなる．
  学習途上で境界点が $\alpha$ の右側（安全域）に移動した場合，その境界点の勾配がゼロとなり，下側裾 $[0, \alpha]$ へ移動しなくなる．
  \item \textbf{一般化逆関数の平坦化（代表点の一意性の喪失）}：
  $g_\alpha(\tau)$ は区間 $[\alpha, 1]$ においてずっと一定値 $1$ をとり続ける（平坦なグラフ）．
  そのため，逆関数 $g_\alpha^{-1}(1)$ は一意な値を返せず，区間 $[\alpha, 1]$ 全体が該当してしまう．
  その結果，式\eqref{eq:gmedian}の $g$-中央値 $\hat\tau_i$ が数学的に定義できなくなる．
\end{enumerate}

\subsubsection{正則化歪み $g_{\alpha,\lambda}$ の設計}
この問題に対し，本研究では一様重みを混合した \textbf{正則化歪み関数} を導入する．
\begin{equation}
  g_{\alpha,\lambda}(\tau) = (1-\lambda)\,g_\alpha(\tau) + \lambda\tau,\qquad \lambda \in (0, 1).
  \label{eq:g-reg}
\end{equation}
導関数を計算すると，
\begin{equation}
  g'_{\alpha,\lambda}(\tau) = \begin{cases} \dfrac{1-\lambda}{\alpha} + \lambda & 0 \le \tau < \alpha \\[0.2cm] \lambda & \alpha < \tau \le 1 \end{cases}
  \label{eq:gprime-reg}
\end{equation}
となる（本実験では $\lambda=0.05$ を既定値とする）．

この正則化により，上記の問題が解消される．
\begin{itemize}
  \item 全域において $g'_{\alpha,\lambda}(\tau) \ge \lambda > 0$ が保証されるため，$\alpha$ の外側にある境界点にも常に $\lambda$ の微小な勾配が供給され続け，勾配消失が完全に回避される．
  \item $g_{\alpha,\lambda}$ は狭義単調増加関数となるため，逆関数 $g_{\alpha,\lambda}^{-1}$ が $[0, 1]$ の全域で一意に定義され，$g$-中央値が常に厳密に計算できる．
\end{itemize}
対応する歪みリスク尺度は $\rho_{g_{\alpha,\lambda}}(Z) = (1-\lambda)\CVaR_\alpha(Z) + \lambda\E[Z]$ であり，CVaR と期待値の凸結合という解釈しやすいリスク評価基準に対応している．

%--------------------------------------------------
\subsection{FQF への実装（歪み FQF）}
\label{subsec:impl}

以上の理論的要素を統合した提案手法（歪み FQF）における1ステップの学習手順を以下に整理する．

\begin{enumerate}
  \item \textbf{quantile fraction の生成}：
  状態観測 $s$ から特徴量 $\psi(s)$ を抽出し，FPN $\omega_1$ に通して $N$ 次元のベクトルを得る．
  ソフトマックス関数を適用して区間幅 $p \in \Delta^{N-1}$ とし，累積和をとって狭義単調増加列 $\tau_0=0 < \tau_1 < \dots < \tau_N=1$ を得る（式\eqref{eq:fpn-cumsum}）．
  \item \textbf{$g$-中央値の算出}：
  各区間 $i=0,\dots,N-1$ に対し，正則化歪み関数を用いて $g$-中央値
  \begin{equation}
    \hat\tau_i = g_{\alpha,\lambda}^{-1}\!\left(\frac{g_{\alpha,\lambda}(\tau_i) + g_{\alpha,\lambda}(\tau_{i+1})}{2}\right)
  \end{equation}
  を計算する（式\eqref{eq:gmedian}）．
  \item \textbf{分位点回帰による $\omega_2$ の更新}：
  Quantile Value Network $\omega_2$ に代表レベル $\hat\tau_i$ を入力し，分位点値 $q_i = \Finv_{\omega_2}(\hat\tau_i; s, a)$ を出力する．
  ターゲット側の離散分布の原子 $q'_j = \Finv_{\omega_2^-}(\hat\tau'_j; s', a^*)$ と区間幅 $w'_j = \tau'_{j+1} - \tau'_j$ に対し，以下の重み付き Huber 分位点回帰損失
  \begin{equation}
    \mathcal{L}_{\mathrm{QR}} = \frac{1}{N} \sum_{i=0}^{N-1} \sum_{j=0}^{N-1} w'_j\,\rho^\kappa_{\hat\tau_i}\!\left(r + \gamma q'_j - q_i\right)
    \label{eq:b-loss}
  \end{equation}
  を計算し，勾配降下法で $\omega_2$ を更新する（$\rho^\kappa_\tau$ は式\eqref{eq:huber-quantile}の分位点 Huber 損失）．ここでターゲット側の重み $w'_j$ は FPN が出力した区間幅であり，離散確率測度としての積分を正確に反映している．
  \item \textbf{代理損失による FPN $\omega_1$ の更新}：
  式\eqref{eq:grad-val}の勾配値 $d_i = g'_{\alpha,\lambda}(\tau_i)[2\Finv(\tau_i) - \Finv(\hat\tau_i) - \Finv(\hat\tau_{i-1})]$ を \texttt{torch.no\_grad()} 下で数値計算する．
  代理損失 $L_{\mathrm{frac}} = \sum_{i=1}^{N-1} \operatorname{stopgrad}(d_i)\,\tau_i$ を作成し，逆伝播を行って FPN のパラメータ $\omega_1$ を更新する．
  \item \textbf{リスク回避的行動選択}：
  環境との相互作用におけるエージェントの行動選択は，純粋な CVaR 歪み $g_\alpha$ を適用した式\eqref{eq:rho-discrete}により，
  \begin{equation}
    a^* = \arg\max_a \sum_{i=0}^{N-1}\left(g_\alpha(\tau_{i+1}) - g_\alpha(\tau_i)\right)\Finv_{\omega_2}(\hat\tau_i; s, a)
  \end{equation}
  として決定論的に選択する．
\end{enumerate}

従来の FQF に対する変更点は，ステップ2の代表点の計算（算術中点から $g$-中央値へ）と，ステップ4の勾配計算（$g'(\tau_i)$ の乗算）の2箇所のみである．
追加の計算量は $O(N)$（$N$ は通常 16 や 32 程度）であり，ネットワーク全体の計算負荷に対して十分に小さい．

%--------------------------------------------------
\subsection{新規性調査の結果要約}
\label{subsec:novelty}

本提案の新規性について，
(a) FQF とリスク尺度の結合，
(b) $\tau$ のサンプリング・サポート点配置の $\Wone$ 最適化，
(c) 歪み重み付き $\Wone$ と損失設計，
(d) 統計・最適輸送における歪み尺度の最適量子化，
の4つの観点から徹底的な学術調査を行った．

調査方法は，arXiv および Google Scholar を主要データベースとし，
\begin{quote}
  \textit{distributional reinforcement learning, risk measure, distortion, quantile optimization, Wasserstein, CVaR, spectral risk, quantile fraction}
\end{quote}
などのキーワードを網羅的に組み合わせて 2026 年 9 月時点までの関連文献を検索した．
選定基準に合致した約 70 文献を精査し，関連度ごとに 10 グループに体系化した．
その結論として，\textbf{FPN の目的関数そのものを歪み重み付き Wasserstein 距離に置換し，quantile fraction学習を歪み分布の最適量子化として定式化した先行研究は，本調査の範囲において確認できなかった}．

近接する代表的先行研究との厳密な差分を表\ref{tab:novelty}に整理する．
\begin{table}[htbp]
\centering
\caption{近接する先行研究との差分一覧．}
\label{tab:novelty}
\small
\setlength{\tabcolsep}{3pt}
\begin{tabular}{>{\raggedright\arraybackslash}p{3.8cm}>{\raggedright\arraybackslash}p{5.0cm}>{\raggedright\arraybackslash}p{5.6cm}}
\toprule
文献 & 関連する内容 & 本提案との差分\\
\midrule
FQF \cite{yang2019fqf} & quantile fractionを一様 $\Wone$ 最小化で学習 & 目的関数が一様（リスク尺度と無関係）．本提案は $g$ で重み付けし歪み最適量子化を達成\\
FQF-SRM \cite{gan2026energy} & FQF とスペクトルリスク尺度（SRM）による行動選択 & FPN の損失は一様 $W_1$ のまま（行動選択のみ歪める）．本提案は FPN 損失そのものを置換\\
IQN \cite{dabney2018iqn} & 行動選択時に $\tau$ のサンプリング分布を歪めてリスク感応方策 & 学習時の $\tau$ 配置は一様（表現力の無駄遣い）．本提案は学習時から表現力を下側裾に最適配分\\
Beyond CVaR \cite{moghimi2025beyond} & 静的スペクトルリスク尺度を最適化する分布型RL & 最適化対象は方策のみ．近似ノードの再配置は行わない\\
Tail-Safe \cite{zhang2025tailsafe} & ヘッジングで $\tau$ サンプリングを裾に集中させる TCC & 発見的な温度制御サンプリング．最適量子化に基づく理論的根拠をもたない\\
EX-DRL \cite{malekzadeh2024exdrl} & GPD（一般化パレート分布）を用いた極値モデリング & 極値理論によるパラメトリックな裾推定．quantile fraction の最適量子化ではない\\
最適離散化理論 \cite{graf2000quantization,pflug2014multistage} & Wasserstein 距離に基づく確率分布の最適量子化 & 静的な確率論・数値解析の枠組み．強化学習の動的な fraction 学習機構へ接続した点が本研究の独自性\\
Lim and Malik \cite{lim2022distributional} & リスク感応方策のための分布型RL（時間不整合の指摘） & 方策最適化の枠組み．表現力そのものの幾何学的再配分ではない\\
RQIQN \cite{zhang2026rqiqn} & quantile geometry の正則化（Wasserstein ロバスト推定） & fraction 依存の補正による頑健化手法．目的関数をリスク尺度に応じて置換する本提案とは機構が異なる\\
Rockafellar and Uryasev \cite{rockafellar2000cvar} & CVaR の変分表現と凸最適化 & 最適化手法の基礎理論．深層ニューラルネットワークによる分布近似誤差の議論は含まない\\
\bottomrule
\end{tabular}
\end{table}

例えば，Gan ら (2026) \cite{gan2026energy} の FQF-SRM は，FQF の出力からスペクトルリスク尺度を計算して行動選択を行っているが，FPN 自体の目的関数は従来通りの一様 $W_1$ のままであり，表現力は一様に散らばったままである．
また，Zhang (2025) \cite{zhang2025tailsafe} の Tail-Safe Hedging や Malekzadeh ら (2024) \cite{malekzadeh2024exdrl} の EX-DRL など，金融工学において裾野を重視する深層強化学習手法は提案されているものの，いずれもサンプリングの工夫や極値統計モデルの外付けであり，\textbf{「quantile fraction の配置を歪み分布の最適量子化として理論づけ，FPN の目的関数そのものを置換した研究」} は存在しない．
Graf and Luschgy \cite{graf2000quantization} の最適量子化理論を FQF のquantile fraction学習と結びつけ，リスク感応型強化学習の表現設計に応用した点に，本研究の独自性がある．
%==================== 5 数値実験 ====================
\section{数値実験}
\label{sec:exp}
本章では，前章までに構築した理論的枠組みとアルゴリズムの有効性を，2つの相補的な数値実験を通じて検証する．
第1の実験（実験A）は，真の確率分布が既知である理想的な設定における \textbf{1次元最適量子化実験} である．
ここでは，強化学習に特有の探索のばらつきやベルマン更新の不安定性を排除し，「quantile fractionの配置基準を一様重みから歪み重み付きへと置換すること」がもつ純粋な数理的近似能力と収束特性を単離して検証する．
第2の実験（実験B）は，自作の連鎖型リスク意思決定環境（RiskyChain-v0）を用いた \textbf{小規模強化学習実験} である．
ここでは，ニューラルネットワークによる関数近似，時間差分学習（TD学習），および探索が組み合わさった動的環境において，提案手法（歪み FQF）が実際に下側裾に境界を適応配置し，CVaR 推定精度と行動選択性能を向上させられるかを検証する．

なお，実験の再現性を担保するため，本研究のすべての計算は AMD Ryzen 7 5825U（8コア16スレッド，32\,GB RAM，CPU演算のみ，PyTorch CPU版）上で実行された．

\subsection{実験A：歪み分布の1次元最適量子化}
\label{subsec:expA}
実験Aの目的は，真の分布関数および分位点関数 $\Finv$ が既知であり，部分期待値が高精度に評価できる環境下において，\textbf{quantile fraction の配置方法の違いだけで $\CVaR_\alpha$ の推定誤差がどれだけ改善されるか} を定量化することである．
深層学習の最適化ノイズやサンプリングの偏りを排し，提案モデルの数理的根拠（命題\ref{prop:bound}の上界および命題\ref{prop:gmedian}の $g$-中央値原子）そのものの有効性を検証する．
したがって実験Aは，真の $\Finv$ と閉形式部分期待値を用いる \textbf{真値参照の歪み最適量子化実験} であり，後述の深層強化学習（実験B）とは性質の異なる純粋な数値解析実験である．

\subsubsection{実験手順}
実験Aは以下の手順で実施した．

\begin{enumerate}
  \item \textbf{対象分布の選定}：
  現実の金融工学や制御工学において現れる代表的なリスク構造を網羅するため，裾の形状や非対称性が大きく異なる以下の3つの確率分布を用意した．
  \begin{itemize}
    \item \textbf{D1：ガウス混合分布} $0.9\,\mathcal N(0,1) + 0.1\,\mathcal N(-4,1)$．
    大部分（90\%）は標準正規分布に従うが，下側裾（$\tau < 0.1$）に平均 $-4$ の小集団が潜んでいる．確率密度が局所的に急変する多峰性をもち，一様な粗い格子ではこの局所的な崖の構造を見落としやすい．
    \item \textbf{D2：負の対数正規分布} $Z = e^{\sigma^2/2} - \exp(\sigma\xi)$，$\xi \sim \mathcal N(0,1)$，$\sigma=1$（平均0）．
    資産収益率のモデル化などで頻出する，左側（損失側）にのみ長く重い裾が非対称に伸びる分布である．
    \item \textbf{D3：$t$分布}（自由度 $\nu=3$）．
    分散は有限（3）であるが4次以上のモーメントが無限大となるヘビーテイル（裾の重い分布）の代表例である．まれに巨大な損失が発生する市場の下側裾リスクを模擬している．
  \end{itemize}

  \item \textbf{参照真値の高精度算出}：
  各分布の真の $\CVaR_\alpha(Z)$ を高精度に算出した．部分期待値 $\mathrm{PE}(a,b) = \E[Z\,\mathbb{I}\{a < Z \le b\}]$ の解析的閉形式と，累積分布関数の数値反転（許容誤差 $10^{-12}$ の高精度求根アルゴリズム）から得た真の分位点 $\Finv(\alpha)$ を用い，$\CVaR_\alpha(Z) = \alpha^{-1}\mathrm{PE}(-\infty, \Finv(\alpha))$ として真値を求めた（ガウス混合 D1 の $\Finv$ 自身は閉形式をもたないため数値求根による）．
  さらに，独立な $10^7$ サンプルのモンテカルロ法（seed 0）を実行し，$|\text{参照値} - \text{MC}| < 3\,\mathrm{SE}$（SE は標準誤差）を満たすことを全設定で確認した．

  \item \textbf{格子点数とリスク水準}：
  原子数（区間数）として $N \in \{8, 16, 32, 64, 128, 256\}$（2の冪乗系列），評価リスク水準として $\alpha \in \{0.01, 0.05, 0.1\}$ の計18通りの組み合わせを設定した．

  \item \textbf{比較手法の構成}：
  quantile fraction $\tau_i$ の配置方針として，以下の4手法を比較した．
  \begin{itemize}
    \item \textbf{M0（一様固定配置）}：QR-DQN 型の等間隔配置 $\tau_i = i/N$ であり，代表原子として算術中点 $\hat\tau_i = (2i+1)/(2N)$ をとる．最適化を行わない基準格子である．
    \item \textbf{M1（$g$-一様配置（$g$-uniform））}：歪み関数の一般化逆関数を用いて $\tau_i = g^{-1}(i/N)$ とする配置である．最適化は行わないが，下側裾に重点的に点を撒くという事前知識のみを与えた対照である．
    \item \textbf{M2（提案法・歪み最適配置）}：本研究の提案法に基づき，目的関数 $L_g$ を最小化するように最適化された配置である．
    単調増加性を保証するため $\tau = \mathrm{cumsum}(\mathrm{softmax}(\theta))$（$\theta \in \R^N$ は最適化変数）と表現し，式\eqref{eq:dist-grad}の境界勾配とソフトマックスのヤコビ行列を用いた連鎖律により，制約付き準ニュートン法（L-BFGS-B法，最大反復500，勾配許容誤差 $10^{-10}$）を実行した．局所解への収縮を防ぐため，M1由来の初期値およびDirichlet摂動による乱数初期値4つの計5つの初期値から最適化を行い（マルチスタート法），最小の損失 $L_g$ を達成した解を採用した．
    \item \textbf{M3（Lloyd型固定点反復）}：最適量子化理論の交互最適化（境界を隣接原子の中点へ更新し，原子を区間の $g$-中央値へ更新する手続きを200回反復）を M2 の独立クロスチェックとして実装した．
  \end{itemize}

  \item \textbf{目的関数の厳密評価}：
  数値積分の台形則近似による誤差の混入を防ぐため，目的関数 $L_g$ は各小区間の積分を部分期待値の閉形式を用いて厳密に計算した．最適化には正則化歪み $g_{\alpha,\lambda}$（$\lambda=0.05$）を用いた．

  \item \textbf{評価指標}：
  評価は正則化を含まない純粋な $g_\alpha$ で行い，絶対誤差 $e = |\CVaR_\alpha(\hat Z_N) - \CVaR_\alpha(Z)|$ および一様固定に対する改善比 $R = e_{\mathrm{M0}} / e_{\mathrm{M2}}$ を算出した．

  \item \textbf{点数換算値（N-equivalence）の評価}：
  提案手法 M2 が $N$ 点で達成した精度と同等以下の誤差を，一様固定配置 M0 が達成するために何点の格子点 $N'$ を必要とするかを求めた．探索範囲 $N' \le 1024$ において二分探索を用いて最小の必要点数 $N'$ を同定した．

  \item \textbf{収束レートの回帰分析}：
  $\alpha=0.05$ において，$N=8$ から拡張測定 $N=1024$ までのデータを用い，$\log e$ を目的変数，$\log N$ を説明変数とする線形回帰（OLS）を行い，誤差減衰の経験的オーダー（直線の傾き）と決定係数 $R^2$ を算出した．

  \item \textbf{有限サンプル設定}：
  真の分位点関数 $\Finv$ が未知で，$n \in \{10^3, 10^4, 10^5\}$ 個の有限データから得られる経験分位関数しか利用できない現実的な状況を想定し，有限サンプル上での M2 の最適化精度を 10 シードで評価した．

  \item \textbf{アブレーション解析}：
  混合比率 $\lambda \in \{0, 0.01, 0.2\}$ の感度分析，および金融工学で広く用いられる連続な歪み尺度である Wang 変換 $g_\eta(\tau) = \Phi(\Phi^{-1}(\tau)+\eta)$（$\eta=1.0$）への適用実験を実施した．
\end{enumerate}

\subsubsection{結果}

\paragraph{分位点配置の視覚化（ラグプロット）}
図\ref{fig:placement}は，D1（ガウス混合），$N=16$（内部境界15点），$\alpha=0.05$ における各手法の quantile fraction 配置を累積確率軸上に短い縦線で示したラグプロット（rug plot）である．

\begin{figure}[htbp]
\centering
\includegraphics[width=0.85\linewidth]{fig_placement.png}
\caption{D1（ガウス混合），$N=16$，$\alpha=0.05$ における各 quantile fraction の配置状況．上から M0（一様固定配置），M1（$g$-一様配置（$g$-uniform）），M2（歪み最適配置）．網掛けは下側裾 $[0,\alpha]$．}
\label{fig:placement}
\end{figure}

一様固定配置 M0（最上段）では，格子幅が $1/16 = 0.0625$ であるため，下側裾 $[0, 0.05]$ の内部には境界点が1点も存在しない．
すなわち，最悪5\%のリスク領域は第1区間 $[0, 0.0625]$ の中に完全に埋没し，下側裾のCVaR評価は単一の中点原子 $\Finv(1/32) = \Finv(0.03125)$ の1点のみで大雑把に代表されている．
これに対し，提案手法 M2（最下段）は，15個の内部境界のうち実に13個を下側裾 $[0, 0.05]$ の内部に集中的に配置している．
残る2個の境界のみが，正則化項 $\lambda=0.05$ による重みに応じて残りの安全領域 $[0.05, 1.0]$ に配分されている．
これは，第\ref{sec:proposed}節で数理的に導出した「提案手法の最適化が，歪み分布 $\dist Z$ に対する最適量子化問題を解いている」という事実を直接的に裏付ける結果である．

\paragraph{CVaR 推定誤差と改善比}
表\ref{tab:A1}に，$\alpha=0.05$，$N=16$ における主結果を示す．
3つの確率分布すべてにおいて，改善比 $R = e_{\mathrm{M0}} / e_{\mathrm{M2}}$ は事前に設定した合格基準 $R \ge 2$ を大幅に超過した．

\begin{table}[htbp]
\centering
\caption{実験A主結果（$\alpha=0.05$，$N=16$）．$e$ は $\CVaR_\alpha$ 推定の絶対誤差，$R=e_{\mathrm{M0}}/e_{\mathrm{M2}}$．M0 は一様固定配置（QR-DQN 型の比較基準）．M1 は裾の事前知識のみを与えた対照（$g$-一様配置，最適化なし）．M2 は真の $\Finv$ を用いる真値参照設定である．判定列 $R\ge2$ は全格子での最低基準である．}
\label{tab:A1}
\small
\setlength{\tabcolsep}{5pt}
\begin{tabular}{lrrrrcc}
\toprule
分布 & $e_{\mathrm{M0}}$ & $e_{\mathrm{M1}}$ & $e_{\mathrm{M2}}$ & $R$ & 判定（$R\ge2$） & M2の $[0,\alpha]$ 内q.f.数\\
\midrule
D1 ガウス混合 & 0.309 & 0.0626 & 0.00316 & 97.9 & PASS & 13 / 15（M0 は 0）\\
D2 負の対数正規 & 2.120 & 0.2450 & 0.04040 & 52.4 & PASS & 13 / 15\\
D3 $t$分布（自由度3） & 0.974 & 0.1480 & 0.02460 & 39.6 & PASS & 13 / 15\\
\bottomrule
\end{tabular}
\end{table}

特にガウス混合分布 D1 においては，改善比は $R = 97.9$（約100倍の精度向上）に達した．
これは，D1 の下側裾 $\tau < 0.1$ に存在する小さな正規分布成分 $\mathcal{N}(-4, 1)$ の崖のような急変構造を，一様固定配置 M0 の単一原子では捉えきれなかったのに対し，提案手法 M2 は13個の小区間を切り分けて局所的な分布形状を忠実に近似できたためである．
また，裾への事前配分のみを行った M1 に対しても，提案法 M2 はすべての分布で誤差を $1/6 \sim 1/20$ に低減しており（$e_{\mathrm{M2}}/e_{\mathrm{M1}} = 0.05 \sim 0.17$），単に裾に点を集めるだけでなく，境界勾配に基づいて各区間幅を最適に微調整することが不可欠であることを裏付けている．
なお，Lloyd 型固定点反復 M3 との照合では，全54設定中46設定で M2 が M3 以下の目的関数値 $L_g$ を達成し，最適化の健全性が確認された．

\paragraph{収束レートとスケーリング則}
図\ref{fig:error}は，点数 $N$ を $8$ から $1024$ まで変化させたときの誤差 $e$ の log-log プロットである．

\begin{figure}[htbp]
\centering
\includegraphics[width=0.95\linewidth]{fig_error.png}
\caption{$\CVaR_\alpha$ 推定誤差 $e$ 対 $N$（log-log）．列はD1/D2/D3，行は$\alpha\in\{0.01,0.05,0.1\}$，線はM0/M1/M2を表す．破線は参照傾き $-1$ である．}
\label{fig:error}
\end{figure}

$\alpha=0.05$ のデータにおいて，両対数グラフの傾き（収束次数）を測定したところ，
一様固定配置 M0 の傾きが $-0.72 \sim -1.27$（おおむね一次精度 $O(N^{-1})$ 程度）にとどまるのに対し，
提案手法 M2 の傾きは $-1.79 \sim -1.84$（決定係数 $R^2 \ge 0.998$）を示し，おおむね二次精度 $O(N^{-2})$ に整合した．
この結果は，注意\ref{rem:midpoint}で変分法により導出した数学的機序と整合している．
すなわち，一様格子では下側裾の端点特異性（$t$分布等の急速な発散）を踏み抜くことで減衰率が $O(N^{-(1-1/\nu)})$（$\nu=3$ で $N^{-2/3}$）へと劣化してしまうのに対し，
歪み最適配置 M2 は第1区間幅を自律的に極小化（$\tau_1 \propto N^{-3}$）して特異性の影響を封じ込めるため，正則部の中点則二次精度 $O(N^{-2})$ を維持できるのである．

\paragraph{点数換算値（N-equivalence）の比較評価}
図\ref{fig:nequiv}は，提案法 M2 が $N$ 点で達成した精度を一様配置 M0 で達成するために何点が必要かを示した N-equivalence の結果である．

\begin{figure}[htbp]
\centering
\includegraphics[width=0.95\linewidth]{fig_nequiv.png}
\caption{点数換算値（N-equivalence，$\alpha=0.05$）．M2 の $N$ 点と同等以下の誤差を一様配置 M0 が得るのに要する点数 $N'$（探索上限は1024点，上限に達した場合は数値を省略）．}
\label{fig:nequiv}
\end{figure}

$\alpha=0.05$ において，D1 ガウス混合では，提案法 M2 のわずか16点と同等の誤差を得るために，一様固定配置 M0 では約695点もの格子点が必要となった．
これは，提案法を用いることで計算リソースやパラメータ数を約43分の1に圧縮できることを意味する．
さらに D2（対数正規）や D3（$t$分布）では，M2 の16点が達成した誤差があまりに小さいため，一様配置は探索上限の1024点を用いても追いつくことができなかった．
この事実は，点数を増やして近似するよりも，目的関数をリスク尺度に合わせて設計することのほうが，同等の精度を達成する上で極めて有効であることを示している．

\paragraph{有限サンプル設定およびアブレーションの示唆}
有限サンプル設定（$n=10^4$）では，10シード平均の改善比の最小値は $R=2.83$ となり，真値参照時に比べて改善幅は縮小した．
これは，サンプル全体のわずか5\%（約500サンプル）しか存在しない下側裾の経験分位関数は統計的な標本分散が大きく，その不安定な領域に多くの原子を配置するためにサンプリングノイズに対する感度が高まるという，表現力と推定量分散のトレードオフを反映している．
それでも大多数の設定で一様配置を上回っており，データから学習する強化学習においても提案法が有効に機能する理論的裏付けとなる．

また，混合比率 $\lambda$ のアブレーションでは，$\lambda=0$（純粋CVaR）で改善比が最大（$R=129.6$）となる反面，分布全体の誤差 $\Wone$ は $3.69$ に悪化した．
これに対し，提案の $\lambda=0.05$ では $R=97.9$ を維持しつつ $\Wone = 0.46$ と分布全体の品質も高く保たれており，下側裾への集中と分布全体の近似のバランスをとる正則化パラメータとしての設計の妥当性が実証された．

\subsection{実験B：小規模強化学習環境における検証}
\label{subsec:expB}
実験Bでは，提案モデルを実際に深層強化学習（FQF）に実装し，ニューラルネットワークによる関数近似，時間差分学習（TD学習），および探索が連動する動的な強化学習環境において，下側裾の $\CVaR_\alpha(s,a)$ 推定精度および獲得方策がどのように改善されるかを検証する．

\subsubsection{評価環境：連鎖型リスク環境（RiskyChain-v0）}
強化学習の検証用として，Gymnasium 形式で連鎖型の意思決定環境（RiskyChain-v0，離散行動）を新規に設計・実装した．

\paragraph{環境の設計思想と破滅リスクの構造}
本環境は，1エピソードの最大ステップ数（ホライズン）$H=10$，エージェントの位置 $x \in \{0, \dots, 9\}$，割引率 $\gamma=1$ の有限ホライズン問題である．
観測状態は，現在位置の10次元 one-hot ベクトルと経過時間割合 $t/H \in [0, 1]$ を連結した11次元ベクトルである．
各ステップにおいて，エージェントは以下の2つの行動から選択を行う．
\begin{itemize}
  \item \textbf{行動 SAFE}：確実に報酬 $0.1$ を獲得し，位置 $x$ は変化しない．
  \item \textbf{行動 RISKY}：確率 $1-p(x)$ で成功して報酬 $0.5$ を獲得し，位置が $x \leftarrow x+1$ と前進する．しかし，位置が進むほど高まる確率 $p(x) = 0.002(x+1)$ でクラッシュ（事故）が発生し，ヘビーテイルに従う極端損失 $-5(1+|T_3|)$（$T_3$ は自由度3の $t$ 分布）を被る（位置は不変）．
\end{itemize}

期待値の観点では，すべての位置において $\E[\mathrm{RISKY}] > 0.1$ となるように設計されている．
したがって，通常のリスク中立エージェント（DQN や平均重視の強化学習）は，目先の高い期待値に釣られて常に RISKY を選び続ける．
しかし，10ステップすべて RISKY を選択した場合，エピソード中に1回以上クラッシュが発生する確率は約 $10.4\%$ に達する．
このクラッシュ確率は，本研究が評価対象とするリスク水準 $\alpha=0.05$（最悪5\%の事態）を大きく上回っている．
したがって，下側5\%の最悪損失の平均である $\CVaR_{0.05}$ を最大化しようとする真のリスク回避的エージェントにとっては，クラッシュを絶対に避けて10ステップすべて SAFE を選択する行動方針（all-SAFE，確定利得 $1.0000$）こそがCVaR最適閾値方策（クラス内最良方策）となる．
「平均的な期待値の高さのみを追求して破滅に至るリスク中立方策」と，「深刻な下側裾リスクを回避して着実に生存を図るリスク回避方策」が対立する点に，本環境の設計意図がある．

\subsubsection{実験手順}
実験Bは以下の設定で実施した．

\begin{enumerate}
  \item \textbf{ネットワーク構造と学習設定}：
  状態次元11を入力とし，隠れ層 $256 \times 256$（ReLU活性化），出力層に $N \times |\mathcal{A}|$ の原子を出力する MLP を用いた．
  原子数は $N=16$，オプティマイザは Adam（主ネットワーク学習率 $10^{-4}$，FPN 学習率 $2.5 \times 10^{-5}$），Huber損失パラメータ $\kappa=1$，リプレイバッファ容量 50k，バッチサイズ 64，ターゲットネットワークのハード更新間隔 500 ステップ，探索率 $\varepsilon$ は $1.0$ から $0.05$ へ 20k ステップで線形減衰させ，総ステップ数は 100k ステップとした．乱数シードは $\{0, 1, 2\}$ の3試行である．
  行動選択およびベルマンターゲットの次状態行動には，$\CVaR_{0.05}$ 推定値を最大化する $\CVaR$-greedy 方策を採用した．

  \item \textbf{ベルマンターゲットの原子重み付け}：
  FQF 系（B0, B1, B1g）においては，ターゲット側の離散分布の原子を区間幅 $w'_j = \tau'_{j+1} - \tau'_j$ で重み付けした Huber 分位点回帰損失（式\eqref{eq:b-loss}）を用いた．
  現在側の代表点には，従来 FQF では中点 $\hat\tau_i$，提案法では $g$-中央値を用いている．

  \item \textbf{比較手法のバリエーション}：
  \begin{itemize}
    \item \textbf{B0（従来FQF）}：一様 $\Wone$ 勾配を用い，FPN を一様配置で初期化．
    \item \textbf{B1（提案歪みFQF・一様初期化）}：提案の歪み境界勾配（式\eqref{eq:dist-grad}）と $g$-中央値原子を採用するが，FPN の初期バイアスは一様配置．
    \item \textbf{B1g（提案歪みFQF・$g$-一様初期化・主結果）}：B1 と同一のモデル構造であるが，FPN の初期バイアスを $g$-一様配置 $\tau_i = g^{-1}_{\alpha,\lambda}(i/N)$ に設定し，学習の初期段階から下側裾に境界を重点配置する．
    \item \textbf{B2（リスク感応IQN）}：標準的な IQN であり，学習時は $\tau \sim U(0, 1)$，行動選択時は $\tau \sim U(0, \alpha)$ からサンプリングする．
    \item \textbf{B2p（IQNサンプリング歪み対抗法）}：B2 の学習時サンプリング自体を $\tau = g^{-1}_{\alpha,\lambda}(U(0, 1))$ と下側裾に偏らせることで，サンプリング側から表現力を集中させようとする素朴な対抗手法．
    \item \textbf{B3 / B3g（固定配置 ablation）}：FPN の重みを凍結し，一様固定配置（B3）または $g$-一様固定配置（B3g）とした対照実験．
  \end{itemize}

  \item \textbf{評価指標の明確な分離}：
  学習の推移を以下の2つの独立した指標で評価した．
  \begin{itemize}
    \item \textbf{M-est（CVaR 推定誤差・主指標）}：エージェントの「頭の中の評価（認知能力）」の正確さを測る指標である．学習済み方策を固定し，全55状態 $\times$ 2行動において $10^4$ 本のロールアウトで高精度に算出した参照値 $\widetilde{\mathrm{CVaR}}_\alpha(s, a)$ と，ネットワークの出力値との平均絶対誤差（MAE）を測定する．
    \item \textbf{M-pol（方策性能・行動結果）}：エージェントの「実際の行動結果」を測る指標である．学習済み方策に従って 5,000 エピソードを実行したときの実際の収益の経験 $\CVaR_{0.05}$ および平均収益を記録する．
  \end{itemize}
\end{enumerate}

\subsubsection{結果}

表\ref{tab:B1}に，100kステップ終了時点における最終性能を示す（3シード平均 $\pm$ 標本SD）．
また，図\ref{fig:B2}に M-est MAE の学習曲線，図\ref{fig:B3}に学習完了後の quantile fraction 配置の分布を示す．

\begin{table}[htbp]
\centering
\caption{実験B最終指標（100k ステップ，3シード平均$\pm$SD，SD は標本 SD（ddof=1））．M-est は $\CVaR_{0.05}$ 推定 MAE（参照値は各手法の greedy 方策下の分布に対するものであることに注意．第\ref{subsec:cross}節参照）．「限定」は閾値方策クラス内最良が切り替わる状態に限定した MAE．}
\label{tab:B1}
\small
\setlength{\tabcolsep}{2.5pt}
\begin{tabular}{lccccc}
\toprule
手法 & M-est（全110） & M-est（限定） & M-pol $\CVaR_{0.05}$ & M-pol 平均 & 下側裾 q.f. 数\\
\midrule
B0 FQF & $7.02\pm0.14$ & $14.41\pm0.15$ & $-10.79\pm0.00$ & $3.75\pm0.00$ & $0.00\pm0.00$\\
B1 歪み FQF（一様初期化） & $1.41\pm0.92$ & $3.24\pm2.28$ & $-1.15\pm2.38$ & $2.03\pm0.65$ & $0.48\pm0.46$\\
\textbf{B1g 歪み FQF（$g$-一様初期化）} & $\mathbf{0.96\pm0.51}$ & $\mathbf{0.39\pm0.29}$ & $\mathbf{0.70\pm0.43}$ & $1.37\pm0.36$ & $\mathbf{14.06\pm0.11}$\\
B2 リスク感応 IQN & $0.36\pm0.05$ & $0.15\pm0.02$ & $1.01\pm0.01$ & $1.13\pm0.22$ & （暗黙的表現）\\
B2p IQN（学習時歪み $\tau$） & $3.42\pm1.38$ & $5.50\pm3.58$ & $1.01\pm0.02$ & $1.13\pm0.22$ & （暗黙的表現）\\
\midrule
B3 FQF（一様固定） & $4.12\pm0.35$ & $9.25\pm0.60$ & $-7.28\pm1.95$ & $3.06\pm0.24$ & $0.00\pm0.00$\\
B3g 歪み FQF（$g$-一様固定） & $0.74\pm0.33$ & $0.66\pm0.52$ & $0.30\pm0.83$ & $1.60\pm0.40$ & $15.00\pm0.00$\\
\midrule
閾値方策クラス内最良（$k^*=0$） & --- & --- & $1.0000$ & $1.0000$ & ---\\
\bottomrule
\end{tabular}
\end{table}

\begin{figure}[htbp]
\centering
\includegraphics[width=0.82\linewidth]{figB2_estimation_ja.png}
\caption{M-estにおける $\CVaR_{0.05}$ 推定 MAE の学習曲線（5手法，3シード平均$\pm$SD，細線は各シード）．}
\label{fig:B2}
\end{figure}

\begin{figure}[htbp]
\centering
\includegraphics[width=0.88\linewidth]{figB3_fractions_ja.png}
\caption{学習済みquantile fraction配置の分布（B0 一様FQF，B1 一様初期化歪みFQF，B1g $g$-一様初期化歪みFQF）．網掛けは下側裾 $[0,\alpha]$．B1g は全シードで理論配置（14〜15個）を忠実に達成している．}
\label{fig:B3}
\end{figure}

\paragraph{B1g の推定精度向上と初期化問題の克服}
主結果である提案法 B1g（$g$-一様初期化）は，全110状態行動対における M-est MAE において $0.96 \pm 0.51$ を記録し，従来 FQF（B0）の $7.02 \pm 0.14$ から約7分の1に誤差を激減させた．
特に，行動選択の命運を分ける意思決定状態（限定領域）においては，B0 の $14.41 \pm 0.15$ に対し，B1g は $0.39 \pm 0.29$ と約37分の1の大幅な推定精度向上を達成した．
学習後の下側裾 $[0, \alpha]$ 内の quantile fraction 数は $14.06 \pm 0.11$ 個（15個中）に達し，3シードすべてで理論配置（約14個）を正確に実現している（図\ref{fig:B3}）．

ここで注目すべき知見は，\textbf{FPN の初期化の役割} である．
一様初期化を採用した旧 B1 では，seed 0 において下側裾 $[0, 0.05]$ に境界が1点も入り込めず，局所解トラップに停滞して学習に失敗した（裾内0個，MAE 2.44，M-pol $-3.81$）．
なぜこの局所解が生じるのかというと，初期状態で境界が安全域 $[0.05, 1]$ に散らばっている場合，下側5\%の狭い領域には初期配置において境界が存在せず，境界勾配に下側裾の損失情報が届かず，勾配消失によって境界が動けなくなる「探索の壁」が存在するためである．
これに対し，初期バイアスを理論配置に合わせる $g$-一様初期化（B1g）を導入したことで，学習の第1ステップから境界が下側裾に配置され，全シードで局所解が完全に解消された．
その結果，実際の方策性能 M-pol $\CVaR_{0.05}$ においても全シードで正値（平均 $0.70 \pm 0.43$）を達成し，危険なクラッシュを回避する健全なリスク回避方策の獲得に成功した．

\paragraph{サンプリング歪み対抗法（B2p）の破綻メカニズム}
一方，比較対象である B2 リスク感応 IQN は M-est MAE $0.36 \pm 0.05$ と高い推定精度を維持した．
これに対し，サンプリング側のみを歪めた対抗手法 B2p は，MAE が $3.42 \pm 1.38$ と通常の B2 から約10倍も悪化し，激しい性能劣化を起こした．

この B2p における推定精度の悪化は，強化学習の学習プロセスを考える上で重要な現象である．
その原因は，\textbf{ベルマンターゲットの生成プロセス} にある．
現在側の分位レベル $\tau_i$ のサンプリング分布を歪めるだけであれば，各レベルにおける Huber 分位点回帰の一致性は崩れない．
しかし，次状態のターゲット側サンプリング $\tau'_j \sim g^{-1}(U)$ を等重みで平均してしまうと，ベルマン更新のターゲット分布そのものが，本来の $r+\gamma Z'$ ではなく，人工的に歪められた架空の分布 $\dist(r+\gamma Z')$ に置き換わってしまう．
その結果，ネットワークは本来の分位点関数 $\Finv$ ではなく，歪み合成された分位点関数 $\Finv(g^{-1}(\tau))$ を学習してしまう．
学習後に評価グリッド $\tau \in [0, 0.05]$ を入力した際，ネットワーク内部では $g^{-1}_{\alpha,\lambda}(\tau) \approx 0.05\tau / 0.95$ という遥かに深い極端損失の分位点を出力してしまい，CVaR を系統的に過小評価して大暴落を起こしたのである．
この結果は，「単にサンプリング頻度を偏らせるだけのアプローチ」はベルマン更新の収束先を狂わせて推定精度を悪化させることを示しており，本研究の提案法のように，目的関数と境界・重みを測度論的に整合させて明示的に制御することの正当性を決定的に証明している．

\paragraph{小規模環境における固定配置アブレーションの分析}
さらに，FPN を凍結した固定配置 ablation（B3, B3g）の結果は，小規模環境特有の性質を浮き彫りにした．
一様固定の B3（MAE 4.12）に対し，理論配置に固定した B3g は MAE $0.74 \pm 0.33$ を達成し，FPN を学習させた B1g（$0.96 \pm 0.51$）と同等以上の推定精度を示した．
これは，本小規模環境（RiskyChain-v0）では all-SAFE 方策下の収益がほぼ確定値 $1.0$（点質量分布）となっており，状態間での収益分布の形状変化が非常に小さいためである．
状態によらずほぼ同一の理論配置が最適である環境では，FPN が状態に応じて境界を動かす余地があまり残されておらず，理論配置を固定するだけで十分な精度が得られてしまう．
後述するように，状態ごとに直面するリスクの裾野が激変する複雑な環境（実験Cのオプションデルタヘッジ）においてこそ，FPN の状態適応性が真の真価を発揮すると結論づけられる．

\subsubsection{共通方策下の交差評価}
\label{subsec:cross}
表\ref{tab:B1}における M-est の結果を解釈する上で，見過ごすことのできない重要な問題が存在する．
それが \textbf{推定対象の難易度に関する交絡（Confounding）} である．

\paragraph{交絡問題の本質}
M-est の参照真値は，各エージェントが獲得した greedy 方策の下での収益分布に対して計算される．
ここで，B0（従来FQF）は CVaR を過大評価した結果，無謀にも all-RISKY 方策を選択している．そのため，B0 が直面している収益分布はクラッシュが頻発する「極めて重い裾をもつ危険な分布」である．
一方，B1g や B2 は賢明にも all-SAFE 方策を選択しているため，直面している収益分布はほぼ「確定値 $1.0$（分散がほぼゼロの平穏な分布）」である．
すなわち，B0 の MAE（7.02）と B1g の MAE（0.96）の差には，手法の表現能力の差だけでなく，「B0 は荒波の中で激しい変動を推定させられているのに対し，B1g は穏やかな海面を推定している」といったような，推定対象そのものの難易度の違いが交絡している．

\paragraph{交差評価のプロトコルと結果}
この交絡を分析するため，エージェントの行動方針を固定した2つの極端な共通参照方策（all-SAFE 方策および all-RISKY 方策）を設定し，
$10^5$ ロールアウト $\times$ 10 独立バッチのモンテカルロ法で算出した高精度参照値と，各学習済みネットワークの出力を全55状態 $\times$ 2行動で比較した（表\ref{tab:B2}）．

\begin{table}[htbp]
\centering
\caption{共通方策下の CVaR 推定 MAE（学習済みネットワーク出力 vs 高精度 MC 参照値，55状態$\times$2行動の平均，3シード平均$\pm$SD，ddof=1）．参照値は $10^5$ ロールアウト $\times$ 10 独立バッチ．各ネットワークは学習時の CVaR-greedy 方策を継続方策として学習しているため，参照方策と一致しない項目には継続方策の不一致による差が混入する（本文参照）．}
\label{tab:B2}
\small
\setlength{\tabcolsep}{4pt}
\begin{tabular}{lcc}
\toprule
評価器（学習済み手法） & 参照方策 all-SAFE（$k=0$） & 参照方策 all-RISKY（$k=10$）\\
\midrule
B0 FQF & $1.88\pm0.14$ & $7.01\pm0.14$\\
B1 歪み FQF（一様初期化） & $0.71\pm0.24$ & $5.78\pm0.30$\\
B1g 歪み FQF（$g$-一様初期化） & $0.59\pm0.10$ & $5.52\pm0.18$\\
B2 リスク感応 IQN & $\mathbf{0.35\pm0.05}$ & $5.31\pm0.12$\\
B2p IQN（学習時歪み $\tau$） & $4.19\pm1.48$ & $4.25\pm0.46$\\
\bottomrule
\end{tabular}
\end{table}

表\ref{tab:B2}から，以下の2つの事実が明らかになった．
\begin{enumerate}
  \item \textbf{推定対象の難易度差が与える影響}：
  比較的平穏な all-SAFE 分布を参照とした場合，B0 の MAE は $7.01$ から $1.88$ へと大幅に低下する．
  逆に，重い裾をもつ all-RISKY 分布を参照とした場合，優秀であったはずの B1g や B2 を含むすべての評価器の MAE が $5.3 \sim 7.0$ へと跳ね上がる．
  この結果は，表\ref{tab:B1}で観測された大きな誤差差の主要因が，重い下側裾そのものの推定難度にあったことを明確に裏付けている．
  \item \textbf{B2p の偏向表現の確認}：
  B2p はいずれの参照方策に対しても約 $4.2$ という大きな誤差を示しており，その学習済み表現が分布全体として歪んでしまっていることが交差評価からも再確認された．
\end{enumerate}

ただし，本交差評価にも留意すべき限界がある．
各ネットワークは学習時の自身の行動方針（継続方策）を前提としてベルマンターゲットを学習しているため，参照方策と学習時の方策が一致しない組み合わせ（例えば all-SAFE を評価する B0）には，方策の不一致による真値のズレが混入する．
交絡を完全に分離するためには，行動選択とベルマン更新の双方を同一方策に固定してゼロから学習し直す「方策評価モード（Policy Evaluation Mode）」での検証が必要であり，これを将来の重要課題として位置づける．

\subsection{考察}
本章の実験Aおよび実験Bの結果から得られた知見を，理論との接続および今後の展開の観点から総括する．

\begin{itemize}
  \item \textbf{数理原理の検証と点数圧縮効果（実験A）}：
  真の分位点関数が既知の設定において，提案法 M2 は一様固定配置 M0 に対して最大約100倍の精度改善を達成した．
  N-equivalence 解析により，提案法の16点は一様格子の約695点（約43倍）に匹敵することが示され，目的関数を歪み尺度に整合させることが計算リソースの大幅な節約につながることが実証された．
  また，収束レートの傾きが一次精度から二次精度 $O(N^{-2})$ へと復帰した現象は，第4章で導出した変分点密度理論（端点特異性の自律的回避）を実証している．

  \item \textbf{強化学習への統合と初期化の必要性（実験B）}：
  強化学習環境 RiskyChain-v0 において，提案法 B1g は従来 FQF（B0）に対して推定誤差を約7分の1（意思決定状態で約37分の1）に低減し，安全な方策の獲得に成功した．
  その過程で，目的関数の変更だけでは初期の勾配消失（局所解トラップ）を招くリスクが明らかになり，理論配置に基づく $g$-一様初期化と組み合わせることで初めて安定した学習が成立するという実践的設計指針が確立された．

  \item \textbf{サンプリング歪み法に対する優位性}：
  IQN のサンプリング頻度を下側裾に偏らせた B2p が約10倍の性能劣化を起こしたのに対し，提案法は安定した高精度を維持した．
  強化学習のベルマン更新において，サンプリング頻度を変更するだけではターゲット分布の固定点が歪んでしまうのに対し，提案法は「原子の配置と重みを測度論的に明示的に制御する」ことで，ベルマン更新の一致性を損なわずに下側裾への表現力集中を両立できることが示された．

  \item \textbf{現状の限界と課題}：
  一方で，小規模環境では収益分布の状態依存性が弱かったため，固定配置 B3g との有意な差は観測されず，FPN の状態適応性の真価は確認されなかった．
  また，連続的な分位レベルを直接入力できる B2 IQN（MAE 0.36）の柔軟性には依然として及んでいない．
  提案法が IQN を凌駕し，状態適応的 FPN の真価を遺憾なく発揮するためには，状態ごとにリスクの裾野が動的に激変する本格的な実応用問題への展開が不可欠である．
\end{itemize}

%==================== 6 今後の展望 ====================
\section{今後の展望}
\label{sec:future}
本研究で確立した「歪み重み付き 1-Wasserstein 距離に基づく quantile fraction 学習」の理論と知見を踏まえ，今後の研究展開として以下の4課題を推進する．

\begin{enumerate}
  \item \textbf{オプションのデルタヘッジ環境における本実験（実験C）}：
  本提案手法の真価を発揮する実応用舞台として，金融工学における \textbf{オプションの動的デルタヘッジ問題} を実装・検証する．
  Black--Scholes モデルおよびその拡張（行使価格 $K=100$，初期原資産価格 $S_0=100$，満期 $T=30$ 営業日，日次時間刻み $\Delta t = 1/252$，ボラティリティ $\sigma=0.2$）のもとで，ATM ヨーロピアン・コールオプション1枚を売り建て（ショート），離散的なヘッジ比率 $a \in \{0, 0.1, \dots, 1.0\}$ で日次リヘッジを行う．
  売買ごとに想定元本の $0.1\%$ の取引コストが発生するため，頻繁な取引はコストを招き，放置は価格変動リスクを招くというトレードオフの下で，終端損益の $\CVaR_{0.05}$ 最大化を目指す．

  \paragraph{なぜデルタヘッジが本手法のキラーアプリケーションなのか}
  実験Bで用いた小規模トイ環境と決定的に異なるのは，\textbf{状態依存性の強さ} である．
  満期までの残り日数 $T-t$ や，原資産価格と行使価格の乖離度 $S_t/K$ に応じて，保有ポジションが直面する損失の確率分布の形状（下側裾の厚さや傾き）は刻一刻と激変する．
  固定格子（B3g）ではこの動的変化に追従できず，状態 $s = (S_t/K, T-t, \text{現ヘッジ比率})$ を入力として quantile fraction $\tau_i(s)$ を伸縮させる FPN（提案手法 B1g）が不可欠となる．
  さらに，理論上完全複製が可能な Black--Scholes 拡散過程にとどまらず，まれに株価が不連続に暴落する Merton のジャンプ拡散過程や Heston 型の確率的ボラティリティモデルを導入することで，現実の金融市場に酷似した極端損失のヘビーテイル環境下で，提案法の優位性を徹底検証する．

  \item \textbf{誤差上界および中点則の数学的定理化}：
  第4章で提示した予想\ref{conj:alpha}（一様格子に対する改善スケーリング比 $\Theta(1/\alpha)$）の数学的証明を完成させ，厳密な定理として確立する．
  本レポートで解明した「一様格子が端点特異性に律速される一方，適応格子は第1区間幅の縮小 $\tau_1 \propto N^{-3}$ によりこれを回避して二次精度 $O(N^{-2})$ を維持する」という変分法による定性的一致（注意\ref{rem:midpoint}）を厳密な不等式評価へと発展させ，符号付き求積誤差の漸近挙動を数学的に証明する．

  \item \textbf{公平な比較対照実験の拡充}：
  実験Bで残された交絡や読み出しプロトコルの差異を完全に解消するため，以下の検証を実施する．
  \begin{itemize}
    \item \textbf{方策評価モードでの交絡分離}：行動選択・ベルマンターゲット・探索のすべてを同一の固定方策に統一して学習を行うことで，推定対象の難易度差を完全に排除した純粋な表現能力比較を行う．
    \item \textbf{読み出しプロトコルの統一}：IQN が評価時に下側裾の等間隔16点を直接読み出している柔軟性に対抗し，FQF の分位点値ネットワークに対しても同様に連続レベルを入力して読み出すプロトコル，および逆に IQN を離散原子風に読み出すプロトコルを比較し，表現形式と読み出し方法の寄与を厳密に切り分ける．
  \end{itemize}

  \item \textbf{リスク尺度の時間整合的な動的強化学習への拡張}：
  各タイムステップで局所的に $\CVaR$-greedy 行動を選択する方策は，エピソード全体の終端損益に対する静的 CVaR 目的とは必ずしも時間整合しない（CVaR の時間不整合性）．
  終端リスクと動的計画法を整合させる拡張状態表現（富の累積値や過去の最大損失を状態に含めるマルコフ決定過程への拡大）と本手法の FPN を統合し，推定精度の改善を真の大域的リスク回避方策の獲得へと直結させる枠組みを構築する．
\end{enumerate}

%==================== 参考文献 ====================
\addcontentsline{toc}{section}{参考文献}
\begin{thebibliography}{99}
\footnotesize
\bibitem{rahaman2019spectral} N. Rahaman, A. Baratin, D. Arpit, F. Draxler, M. Lin, F. A. Hamprecht, Y. Bengio, A. Courville.
On the Spectral Bias of Neural Networks. \textit{ICML 2019}.
\bibitem{yang2019fqf} D. Yang, L. Zhao, Z. Lin, T. Qin, J. Bian, T.-Y. Liu.
Fully Parameterized Quantile Function for Distributional Reinforcement Learning. \textit{NeurIPS 2019}.
\bibitem{dabney2018iqn} W. Dabney, G. Ostrovski, D. Silver, R. Munos.
Implicit Quantile Networks for Distributional Reinforcement Learning. \textit{ICML 2018}.
\bibitem{dabney2018qrdqn} W. Dabney, M. Rowland, M. G. Bellemare, R. Munos.
Distributional Reinforcement Learning with Quantile Regression. \textit{AAAI 2018}.
\bibitem{bellemare2017c51} M. G. Bellemare, W. Dabney, R. Munos.
A Distributional Perspective on Reinforcement Learning. \textit{ICML 2017}.
\bibitem{moghimi2025beyond} M. Moghimi, H. Ku.
Beyond CVaR: Leveraging Static Spectral Risk Measures for Enhanced Decision-Making in Distributional Reinforcement Learning. \textit{ICML 2025}.
\bibitem{zhang2025tailsafe} J. Zhang.
Tail-Safe Hedging: Explainable Risk-Sensitive Reinforcement Learning with a White-Box CBF--QP Safety Layer in Arbitrage-Free Markets.
arXiv:2510.04555, 2025.（単著・査読なしプレプリント）
\bibitem{graf2000quantization} S. Graf, H. Luschgy.
Foundations of Quantization for Probability Distributions. \textit{Lecture Notes in Mathematics} 1730, Springer, 2000.
\bibitem{rockafellar2000cvar} R. T. Rockafellar, S. Uryasev.
Optimization of Conditional Value-at-Risk. \textit{Journal of Risk} 2(3): 21--41, 2000.
\bibitem{lim2022distributional} S. H. Lim, I. Malik.
Distributional Reinforcement Learning for Risk-Sensitive Policies. \textit{NeurIPS 2022}.
\bibitem{yaari1987dual} M. E. Yaari.
The Dual Theory of Choice under Risk. \textit{Econometrica} 55(1): 95--115, 1987.
\bibitem{pichler2013evaluations} A. Pichler.
Evaluations of Risk Measures for Different Probability Measures. \textit{SIAM Journal on Optimization} 23(1): 530--551, 2013.
\bibitem{kiesel2016wasserstein} R. Kiesel, R. R{\"u}hlicke, G. Stahl, J. Zheng.
The Wasserstein Metric and Robustness in Risk Management. \textit{Risks} 4(3): 32, 2016.
\bibitem{gan2026energy} G. Gan, N. Chen, Z. Wei, H. Qin, Y. Lu, H. Huang, Y. Huang.
Energy-aware job scheduling for green data centers under uncertainty: A structure-compressed spectral risk control approach. \textit{Applied Energy} 420: 128137, 2026.
\bibitem{pflug2014multistage} G. C. Pflug, A. Pichler.
\textit{Multistage Stochastic Optimization}. Springer Series in Operations Research and Financial Engineering, Springer, 2014.
\bibitem{cao2021deep} J. Cao, J. Chen, J. Hull, Z. Poulos.
Deep Hedging of Derivatives Using Reinforcement Learning. \textit{The Journal of Financial Data Science} 3(1): 10--27, 2021.
\bibitem{cao2023gamma} J. Cao, J. Chen, S. Farghadani, J. Hull, Z. Poulos, Z. Wang, J. Yuan.
Gamma and vega hedging using deep distributional reinforcement learning. \textit{Frontiers in Artificial Intelligence} 6: 1129370, 2023.
\bibitem{malekzadeh2024exdrl} P. Malekzadeh, Z. Poulos, J. Chen, Z. Wang, K. N. Plataniotis.
EX-DRL: Hedging Against Heavy Losses with EXtreme Distributional Reinforcement Learning. arXiv:2408.12446, 2024.（プレプリント）
\bibitem{zhang2026rqiqn} Z. Zhang, M. Yang, R. Chen, S. Xie, H. Xiong.
Quantile Geometry Regularization for Distributional Reinforcement Learning. arXiv:2605.08182, 2026.（プレプリント）
\end{thebibliography}
\end{document}